\documentclass[11pt,letterpaper]{article}
\usepackage[margin=1in]{geometry}
\usepackage{amsmath,amssymb,amsthm,mathtools}
\usepackage{booktabs,array}
\usepackage[hidelinks]{hyperref}
\usepackage{microtype}
\numberwithin{equation}{section}
\newtheorem{theorem}{Theorem}[section]
\newtheorem{proposition}[theorem]{Proposition}
\newtheorem{lemma}[theorem]{Lemma}
\newtheorem{corollary}[theorem]{Corollary}
\theoremstyle{definition}
\newtheorem{definition}[theorem]{Definition}
\theoremstyle{remark}
\newtheorem{remark}[theorem]{Remark}
\newcommand{\R}{\mathbb R}
\newcommand{\N}{\mathbb N}
\newcommand{\C}{\mathbb C}
\newcommand{\ind}{\mathbf 1}
\newcommand{\dd}{\,\mathrm d}
\newcommand{\Rea}{\operatorname{Re}}
\newcommand{\lcm}{\operatorname{lcm}}
\newcommand{\norm}[1]{\left\lVert#1\right\rVert}
\newcommand{\abs}[1]{\left\lvert#1\right\rvert}
\newcommand{\Lap}{\mathcal L}
\DeclareMathOperator{\diag}{diag}
\DeclareMathOperator{\supp}{supp}
\setcounter{tocdepth}{1}
\setlength{\emergencystretch}{2em}
\hypersetup{pdftitle={Selberg's Formula as a Riccati Flow, with an Application to a Sparse Mertens Certificate},
  pdfauthor={Jeffery Kline},pdfsubject={Version 0.1.0, doi:10.5281/zenodo.23301917}}

\title{Selberg's Formula as a Riccati Flow,\\with an Application to a Sparse
Mertens Certificate}
\author{Jeffery Kline}
\date{October 2026\\[3pt]\small Version 0.1.0, \href{https://doi.org/10.5281/zenodo.23301917}{doi:10.5281/zenodo.23301917}}

\begin{document}
\maketitle
\begin{abstract}
Wirsing showed in 1962--64, by elementary methods built on Selberg's
symmetry formula, that $\psi(x)=x+O(x/\log x)$; Bombieri and Wirsing
reached $O(x/(\log x)^A)$ for every $A$. We give a different proof of the
mean-square form of the first bound. Let $f(t)=e^{-t}\psi(e^t)-1$, and let $b(s)$ be the Laplace
transform of the prime-power measure. On each horizontal line
$s=\sigma+i\xi$, Selberg's formula becomes a Riccati equation for
$\sigma b(\sigma+i\xi)$ as $\sigma$ decreases. Positivity at the doubled
frequency $2\xi$ keeps the solution away from the unstable fixed point $-1$.
A M\"obius substitution then turns it into a linear equation whose forcing
term depends on the solution but stays small, which gives a resolvent bound
uniform in $\xi$. Combined with a local density bound for
primes from a quadratic sieve, this proves the energy bound
$\int_0^T t^2f(t)^2\,\mathrm dt=O(T)$. The proof stays in the half-plane of
absolute convergence and uses no analytic continuation.

The argument applies to the general coupled-source setting of
Theorem~\ref{thm:energy}: a positive measure $\nu$ with $\nu([0,t])=t+O(1)$,
coupled to a bounded source $f$ by a first-order equation and a Selberg-type
relation with bounded error. With a local density hypothesis it gives the
energy bound; without it, the same method proves $f\in L^2(0,\infty)$. This
supplies the hypotheses of Wirsing's abstract Tauberian theorem that the
Selberg relation does not give directly. That theorem, which Wirsing stated
without a detailed proof, then gives $f(t)=O(1/t)$, with a constant that may
depend on the measure. As an application, the energy bound closes a sparse
divisor--cofactor certificate for the Mertens function and gives
$M(N)=O(N/\log N)$. These rates are classical. The contribution is the
method.
\end{abstract}

\noindent\textbf{Status and AI assistance.}
The proofs in this paper were written with the assistance of AI models and
were checked by separate fresh-context AI proof audits, all using models
from one AI provider. They have not been reviewed by a
human expert or formally verified. Novelty is qualified in
Section~\ref{sec:literature}.

\tableofcontents

\section{Introduction}\label{sec:overview}

Let $\mu(n)$ be the M\"obius function, $\Lambda(n)$ the von Mangoldt function,
and
\[
 M(x)=\sum_{n\le x}\mu(n),\qquad
 \psi(x)=\sum_{n\le x}\Lambda(n),\qquad
 \vartheta(x)=\sum_{p\le x}\log p.
\]
All logarithms are natural. A letter $p$ or $q$ used as an arithmetic index
denotes a prime. We reserve $\nu$ for positive measures, so that the measure
in the analytic argument is not confused with $\mu(n)$. Dirichlet convolution
is denoted by $*_D$; convolution on the additive half-line or a specified
circle is denoted by $*$. Constants implicit in $O$, $\ll$, and $\asymp$ are
independent of the scale tending to infinity. They depend only on the
stated hypothesis bounds, except in Corollary~\ref{cor:wirsing-bootstrap},
where they may also depend on the particular pair $(f,\nu)$.

The main analytic result is Theorem~\ref{thm:energy}. In the arithmetic
case it applies to
\begin{equation}\label{eq:arith-source}
 f(t)=e^{-t}\psi(e^t)-1,\qquad
 \nu=\sum_{n\ge2}\frac{\Lambda(n)}n\delta_{\log n},\qquad t\ge0,
\end{equation}
and gives
\begin{equation}\label{eq:main-energy}
 E(T):=\int_0^T t^2|f(t)|^2\dd t\ll T\qquad(T\ge1).
\end{equation}
Equivalently, $\int_T^{2T}|f(t)|^2\dd t\ll T^{-1}$. This is a mean-square
bound. The pointwise bound $f(t)=O(1/t)$, which is stronger, is due to
Wirsing \cite{wirsing1,wirsing2}.

\paragraph{How the proof works.}
Let $b(s)$ be the Laplace transform of $\nu$. On a horizontal line
$s=\sigma+i\xi$, Selberg's formula becomes the Riccati equation
$\beta'=-\beta-\beta^2+e$ for $\beta=\sigma b(\sigma+i\xi)$, where the
derivative is taken as $\sigma$ decreases and the error $e$ is small
(Section~\ref{sec:scale}). The point $\beta=-1$ is an unstable fixed point.
If $\beta$ came close to it, positivity of $\nu$ would force
$\Rea\beta(\sigma,2\xi)$ to stay near $1$. The same Riccati equation at
frequency $2\xi$ drives it down too fast for that to last. So
$\beta$ stays away from $-1$, uniformly in $\xi$
(Lemma~\ref{lem:exclusion}). The substitution $W=\beta/(1+\beta)$ cancels the
quadratic term exactly, leaving a linear equation for $W$ whose forcing
depends on $\beta$ but is small. Solving it bounds
$\int|b(u+i\xi)|\dd u$ along each line, which gives a resolvent for $H$,
the Laplace transform of $f$ (Section~\ref{sec:resolvent}). The remaining
work is to control $H$ on a starting line and at high frequencies
(Sections~\ref{sec:tails}--\ref{sec:closure}). This is where the density
hypothesis \eqref{eq:LD} is used. For primes it follows from a quadratic
sieve (Section~\ref{sec:sieve}).

\paragraph{Without the density hypothesis.}
Section~\ref{sec:wirsing-bridge} drops \eqref{eq:LD}. The same argument then
proves $f\in L^2(0,\infty)$. For a general positive measure, this supplies
the hypotheses of Wirsing's abstract theorem (his Satz~2) that the Selberg
relation does not give directly. That theorem then gives $f(t)=O(1/t)$ and
\eqref{eq:main-energy}, with constants that may depend on the pair
$(f,\nu)$. This step relies on Wirsing's theorem, which he stated without a detailed proof.
Nothing else in the paper depends on it.

\paragraph{The application.}
Sections~\ref{sec:certificate}--\ref{sec:transfer} turn the energy bound
into an exact finite identity with prime and semiprime coefficients. They
handle discrete sampling, coprimality, and the finite cutoff explicitly. The
results are
\begin{equation}\label{eq:arith-conclusions}
 |M(N)|\ll\frac{N}{\log N},\qquad
 \psi(x)=x+O\!\left(\frac{x}{(\log x)^{1/3}}\right).
\end{equation}
These rates are classical.

\paragraph{What we think is new.}
The bounds are not new; the way they are proved is. We see three new
ingredients.
\begin{enumerate}
 \item \emph{Selberg's formula as a differential equation in $\sigma$.}
       The formula is solved as a quantitative Riccati equation for the
       Laplace transform, inside the half-plane of absolute convergence.
       The elementary remainder proofs we examined (Kuhn, Wirsing, Bombieri)
       work with real-variable convolution inequalities; Bombieri uses the
       Laplace transform only to prove a convolution identity on the real
       axis. The analytic proofs we examined (Wiener--Geller, the 3-4-1
       argument, Tao) use the transform on the boundary line and give no
       rate (Section~\ref{sec:lit-mechanisms}). Our argument uses the transform but never reaches the
       boundary, and it gives a rate.
 \item \emph{Positivity at $2\xi$ as a barrier for the flow.} The
       inequality is the one behind the 3-4-1 argument. What is new is how
       it is used. It keeps the flow a fixed distance from $-1$ for a fixed
       time, uniformly in $\xi$, and this replaces analyticity.
 \item \emph{The linearized resolvent and what it gives.} The substitution
       $W=\beta/(1+\beta)$ removes the quadratic term of the flow; the
       remaining forcing depends on the solution but is small, and the
       resulting resolvent is bounded uniformly in $\xi$. In the general
       coupled-source setting it gives $f\in L^2$ with no density hypothesis. Through Wirsing's
       abstract theorem, this makes his pointwise bound $f(t)=O(1/t)$
       available beyond the arithmetic case (Section~\ref{sec:wirsing-bridge}).
\end{enumerate}
We did not find any of these in the sources compared in
Section~\ref{sec:lit-mechanisms}. That comparison was not exhaustive, so
priority is not settled.

\paragraph{What the proof uses.}
Laplace transforms in their half-plane of absolute convergence, Plancherel's
theorem, and the Poisson kernel. It uses no analytic continuation and no
zero-free region, no Hahn--Banach theorem, and no imported bound for the
error in the prime number theorem. All arithmetic inputs are proved or cited
explicitly.

\paragraph{What is not proved.}
The resolvent bound holds for a source built from the same measure as its
convolution multiplier. It is not a bound for arbitrary vectors.
Sections~\ref{sec:finite} and~\ref{sec:recursion} record the finite
eigenstructure, singular-vector formulation, and recursions that motivated
the work, and state which assertions remain open. A related estimate on
the circle, which does not by itself close the argument, is in
Appendix~\ref{app:circle}.

\section{The abstract weighted energy theorem}\label{sec:abstract}

\begin{theorem}[Weighted energy for a coupled positive source]\label{thm:energy}
Let $\nu$ be a positive, locally finite measure on $[0,\infty)$ with no atom
at zero, and set $B(t)=\nu([0,t])$. Suppose $f$ is real, bounded,
right-continuous, and locally of bounded variation, with $f(0)=-1$. Assume
\begin{align}
 B(t)&=t+O(1),\label{eq:B}\\
 \nu&=df+(1+f)\dd t\quad\hbox{on }(0,\infty),\label{eq:source}\\
 R(t)&:=tf(t)+\int_{[0,t]}f(t-u)\nu(\mathrm du)=O(1).\label{eq:R}
\end{align}
Assume further that, for $t\ge0$ and $0<h\le1$,
\begin{equation}\label{eq:LD}
 \nu((t,t+h])\le C_\ell h+q(t),\qquad
 0\le q(t)\le C_q(1+t)^2e^{-t/2}.
\end{equation}
Then $E(T)\ll T$ for all $T\ge1$, with a constant depending only on the
hypothesis bounds.
\end{theorem}

The derivative in \eqref{eq:source} is the Stieltjes derivative on the positive
half-line. When a function is extended by zero to negative arguments, its
boundary value at zero is accounted for separately. In particular,
integrating \eqref{eq:source} gives the bounded primitive
\begin{equation}\label{eq:primitive}
 F(t):=\int_0^t f(u)\dd u=B(t)-t-f(t)-1=O(1).
\end{equation}
Right-continuity fixes representatives at atoms; all subsequent $L^2$
statements are independent of changes on null sets.

For $s=\sigma+i\xi$, $\sigma>0$, define
\begin{equation}\label{eq:transforms}
 b(s)=\int_{[0,\infty)}e^{-st}\nu(\mathrm dt),\quad
 H(s)=\int_0^\infty e^{-st}f(t)\dd t,\quad
 P(s)=\int_0^\infty e^{-st}R(t)\dd t.
\end{equation}
These integrals and their derivatives converge absolutely where used. Our
Fourier transform is $\widehat u(\xi)=\int_\R e^{-i\xi t}u(t)\dd t$, so
$\norm{\widehat u}_2^2=2\pi\norm u_2^2$.

\section{The scale equation and two-frequency positivity}\label{sec:scale}

\subsection{The exact Riccati equation}
Integration by parts, with the boundary value $f(0)=-1$, gives
\begin{equation}\label{eq:BH}
 b(s)=(s+1)\left(H(s)+\frac1s\right),\qquad
 P(s)=-H'(s)+b(s)H(s).
\end{equation}
For example, $\int_{(0,\infty)}e^{-st}df(t)=1+sH(s)$; the constant $1$
must not be dropped. Eliminating $H$ gives
\begin{equation}\label{eq:riccati-b}
 -b'(s)+b(s)^2=A(s)b(s)+D(s)+(s+1)P(s),
\end{equation}
where
\[
 A(s)=1+\frac1{s(s+1)},\qquad D(s)=\frac{s+1}{s^2}.
\]
Fix $\xi\ne0$ and $a>0$, write $\sigma=a e^{-v}$, and put
$\beta(\sigma,\xi)=\sigma b(\sigma+i\xi)$. Equation~\eqref{eq:riccati-b}
becomes
\begin{equation}\label{eq:riccati-beta}
 \frac{\mathrm d\beta}{\mathrm dv}=-\beta-\beta^2+e(\sigma,\xi),\qquad
 e=\sigma A(s)\beta+\sigma^2D(s)+\sigma^2(s+1)P(s).
\end{equation}
Let $g(\xi)=1+|\xi|+|\xi|^{-1}$. There is a constant $C_e\ge1$ such that
\begin{equation}\label{eq:defect}
 |e(\sigma,\xi)|\le C_e\sigma g(\xi),
 \qquad 0<\sigma\le\min(1,|\xi|).
\end{equation}
Indeed $|P(s)|\le\norm R_\infty/\sigma$,
$|A(s)|\le1+|\xi|^{-1}$, and
$|D(s)|\le(2+|\xi|)/|\xi|^2$. Positivity and \eqref{eq:B} give
\begin{equation}\label{eq:csigma}
 |\beta(\sigma,\xi)|\le c_\sigma:=\sigma b(\sigma),\qquad
 |c_\sigma-1|\le C_B\sigma,
\end{equation}
where $C_B\ge1$ bounds $|B(t)-t|$. To see the last estimate, use
$c_\sigma=\sigma^2\int_0^\infty e^{-\sigma t}B(t)\dd t$.
The term $\sigma^2/|\xi|^2$ in the defect is at most $\sigma/|\xi|$ in
the stated range. The high-frequency loss in \eqref{eq:defect} is retained.

\subsection{Uniform exclusion of the multiplier \texorpdfstring{$-1$}{-1}}
Fix, once and for all,
\begin{equation}\label{eq:c-choice}
 0<c\le\min\left\{10^{-3},\frac1{1000C_B},\frac1{2000C_e}\right\},
 \qquad a(\xi)=\frac{c}{g(\xi)}.
\end{equation}
For $0<\sigma\le a(\xi)$, both $\xi$ and $2\xi$ lie in the range of
\eqref{eq:defect}, and
\begin{equation}\label{eq:small-defects}
 |c_\sigma-1|\le10^{-3},\qquad
 |e(\sigma,\xi)|\le10^{-3},\qquad
 |e(\sigma,2\xi)|\le10^{-3}.
\end{equation}
Here $g(2\xi)\le2g(\xi)$ supplies the second-frequency bound. No condition
$\sigma\le a(2\xi)$ is needed.

The measure $e^{-\sigma t}\nu(\mathrm dt)/b(\sigma)$ is a probability measure.
For $U=e^{-i\xi t}$, the identity $\Rea U^2=2(\Rea U)^2-1$ and the
nonnegativity of variance give
\begin{equation}\label{eq:two-freq}
 \Rea\beta(\sigma,2\xi)
 \ge\frac{2(\Rea\beta(\sigma,\xi))^2}{c_\sigma}-c_\sigma.
\end{equation}

\begin{lemma}\label{lem:exclusion}
For every $\xi\ne0$ and $0<\sigma\le a(\xi)$,
$|1+\beta(\sigma,\xi)|\ge1/100$.
\end{lemma}
\begin{proof}
Suppose the distance is less than $1/100$ at a starting value $\sigma_*$.
Follow \eqref{eq:riccati-beta} for $\sigma=\sigma_*e^{-v}$,
$0\le v\le\log2$. For $y=1+\beta(\sigma,\xi)$,
$y'=-\beta y+e$. Since $|\beta|\le1.001<2$, Gronwall's inequality yields
\[
 |y(v)|\le e^{2v}|y(0)|+\frac{e^{2v}-1}{2000}
 <\frac4{100}+\frac3{2000}<\frac1{20}.
\]
Hence $\Rea\beta(\sigma,\xi)<-19/20$. Equation~\eqref{eq:two-freq} gives
$u(v):=\Rea\beta(\sigma,2\xi)\ge4/5$ throughout the interval. The same
scale equation at $2\xi$ and \eqref{eq:small-defects} imply
\[
 u'\le-u-2u^2+|\beta(\sigma,2\xi)|^2+10^{-3}
 \le-\frac45-2\left(\frac45\right)^2+(1.001)^2+10^{-3}<-1.
\]
This is impossible: $u$ must remain in $[0.8,1.001]$, an interval of length
$0.201$, for time $\log2>0.69$.
\end{proof}

\section{A uniform resolvent along the actual source}\label{sec:resolvent}

For a fixed nonzero $\xi$, set $a=a(\xi)$, $\sigma=ae^{-v}$, and
$W=\beta/(1+\beta)$. The exact cancellation of the quadratic term gives
\begin{equation}\label{eq:W}
 W'+W=\frac{e}{(1+\beta)^2}.
\end{equation}
Lemma~\ref{lem:exclusion} bounds the denominator away from zero.
Also $|W(0)|\le101$ and $|e|\le C_e c e^{-v}$, so
\begin{align}
 |W(v)|&\ll(1+v)e^{-v},\notag\\
 |\beta(\sigma,\xi)|&\ll\frac\sigma a\left(1+\log\frac a\sigma\right),
 \label{eq:beta-decay}\\
 |b(\sigma+i\xi)|&\ll\frac1a\left(1+\log\frac a\sigma\right).
 \label{eq:b-decay}
\end{align}
For the second line use $\beta=(1+\beta)W$ and $|1+\beta|\le2.001$.
In particular,
\begin{equation}\label{eq:integrated-b}
 \int_0^{a(\xi)}|b(u+i\xi)|\dd u\ll1
\end{equation}
uniformly in $\xi\ne0$.

Solving the actual source equation $H'=bH-P$ along the horizontal line
gives
\begin{align}
 H(\sigma+i\xi)
 ={}&\exp\!\left(-\int_\sigma^a b(v+i\xi)\dd v\right)H(a+i\xi)
 \notag\\
 &+\int_\sigma^a
  \exp\!\left(-\int_\sigma^u b(v+i\xi)\dd v\right)P(u+i\xi)\dd u.
 \label{eq:source-formula}
\end{align}
Differentiation in $\sigma$ verifies both the sign and integration direction.
Writing $H_*(\xi)=H(a(\xi)+i\xi)$, \eqref{eq:integrated-b} implies
\begin{equation}\label{eq:source-bound}
 |H(\sigma+i\xi)|\ll |H_*(\xi)|+
              \int_\sigma^{a(\xi)}|P(u+i\xi)|\dd u.
\end{equation}
The initial-source term $H_*$ is retained. Its control is an essential part
of the proof, not an independently assumed smallness condition.

\section{Translation estimates and weighted Fourier tails}\label{sec:tails}

\begin{lemma}\label{lem:translations}
For $k\in\{0,1\}$, put $U_{\sigma,k}(t)=t^ke^{-\sigma t}f(t)$ on
$[0,\infty)$ and zero on $(-\infty,0)$. For $0<\sigma\le1$ and $0<h\le1$,
\begin{equation}\label{eq:translation}
 \norm{U_{\sigma,k}(\cdot-h)-U_{\sigma,k}}_2^2
 \ll h^2\sigma^{-2k-1}+h.
\end{equation}
\end{lemma}
\begin{proof}
Solving \eqref{eq:source} over $(t,t+h]$ gives
\[
 f(t+h)-f(t)=(e^{-h}-1)(1+f(t))
       +\int_{(t,t+h]}e^{u-t-h}\nu(\mathrm du).
\]
Consequently $|f(t+h)-f(t)|$ is bounded both by
$Ch+\nu((t,t+h])$ and by $Ch+q(t)$. Multiply these two bounds.
With $w(t)=(1+t)^{2k}e^{-2\sigma t}$, the resulting four terms are controlled
by
\begin{align*}
 h^2\int_0^\infty w(t)\dd t&\ll h^2\sigma^{-2k-1},\\
 h\int_0^\infty w(t)\nu((t,t+h])\dd t&\ll h^2\sigma^{-2k-1},\\
 h\int_0^\infty w(t)q(t)\dd t&\ll h,\\
 \int_0^\infty w(t)q(t)\nu((t,t+h])\dd t&\ll h.
\end{align*}
For the second line, Tonelli's theorem restricts $t$ to an interval of
length at most $h$ adjacent to the integration variable $u$. On that
interval $w(t)\ll(1+u)^{2k}e^{-2\sigma u}$. The measure moment is
$O(\sigma^{-2k-1})$ by $B(u)=u+O(1)$. For the fourth line use the
majorant $(1+u)^{2k+2}e^{-u/2}$ for $w(t)q(t)$ on the same interval;
its integral against $\nu$ is finite uniformly in $\sigma$. This uses the
majorant of $q$, not any unasserted regularity of $q$ itself.

For the change of weight $v_k(t)=t^ke^{-\sigma t}$, Cauchy--Schwarz gives
$\int_0^\infty|v_k(t+h)-v_k(t)|^2\dd t\le h^2\norm{v_k'}_2^2$.
The latter is $O(h^2\sigma)$ for $k=0$ and $O(h^2/\sigma)$ for $k=1$,
both within the stated bound. Finally the boundary strip at zero costs
$O(h^{2k+1})$. These estimates prove \eqref{eq:translation}.
\end{proof}

Average the Fourier translation identity over $0<h<2/K$, $K\ge2$.
For $|\xi|\ge K$ the averaged multiplier is
\[
 2\left(1-\frac{\sin(2\xi/K)}{2\xi/K}\right)\ge1.
\]
Since the transforms of $U_{\sigma,0}$ and $U_{\sigma,1}$ are respectively
$H(\sigma+i\xi)$ and $-H'(\sigma+i\xi)$, we obtain
\begin{align}
 \int_{|\xi|\ge K}|H(\sigma+i\xi)|^2\dd\xi
   &\ll\frac1{\sigma K^2}+\frac1K,\label{eq:tail-H}\\
 \int_{|\xi|\ge K}|H'(\sigma+i\xi)|^2\dd\xi
   &\ll\frac1{\sigma^3K^2}+\frac1K.\label{eq:tail-Hprime}
\end{align}
The terms $1/K$ explicitly retain the atomic and initial-boundary costs.

\section{The variable starting line}\label{sec:start}

\begin{lemma}\label{lem:start}
The initial-source term satisfies
\begin{align}
 \int_{r\le|\xi|\le2r}|H_*(\xi)|^2\dd\xi&\ll r&& (0<r\le1),
 \label{eq:low-start}\\
 \int_{Q\le|\xi|\le2Q}|H_*(\xi)|^2\dd\xi&\ll Q^{-1}&& (Q\ge1).
 \label{eq:high-start}
\end{align}
\end{lemma}
\begin{proof}
By \eqref{eq:primitive},
$H(s)=s\int_0^\infty e^{-st}F(t)\dd t$, whence
$|H(s)|\ll|s|/\Rea s$. On a low shell, $a(\xi)\asymp cr$ and
$|a(\xi)+i\xi|\asymp r$. This proves \eqref{eq:low-start}.

For a positive high shell $Q\le\xi\le2Q$, $Q\ge1$, we have
\[
 \frac c{4Q}\le a(\xi)\le\frac cQ.
\]
Set $\tau=c/(8Q)$ and $d(\xi)=a(\xi)-\tau$. Then
$d_1:=c/(8Q)\le d(\xi)\le7c/(8Q)=:d_2$.
For the Poisson kernel $\mathcal P_d(y)=d/[\pi(d^2+y^2)]$,
\begin{equation}\label{eq:poisson}
 H(a(\xi)+i\xi)=\mathcal P_{d(\xi)}*[H(\tau+i\cdot)](\xi).
\end{equation}
For fixed $d$ this follows by Fourier transformation: convolution by
$\mathcal P_d$ multiplies the original half-line function by
$e^{-d|t|}=e^{-dt}$. Both sides are continuous on the positive half-plane,
so they can be evaluated at each indicated $(d(\xi),\xi)$.

Jensen's inequality bounds the square in \eqref{eq:poisson} by the
Poisson average of $|H(\tau+i\cdot)|^2$. The dependence of $d$ on $\xi$
does not invalidate integration, because
\begin{equation}\label{eq:poisson-ratio}
 \mathcal P_{d(\xi)}(y)\le\frac{d_2}{d_1}\mathcal P_{d_1}(y)
 \le7\mathcal P_{d_1}(y).
\end{equation}
Thus frequencies $\eta\in[Q/2,3Q]$ contribute at most
\[
 C\int_{|\eta|\ge Q/2}|H(\tau+i\eta)|^2\dd\eta\ll Q^{-1}
\]
by \eqref{eq:tail-H} when $Q\ge4$. Outside this enlarged shell,
$|\xi-\eta|\ge Q/2$, so the kernel is $O(Q^{-3})$ and its integral over
$\xi\in[Q,2Q]$ is $O(Q^{-2})$. That contribution is at most
\[
 \frac C{Q^2}\int_\R|H(\tau+i\eta)|^2\dd\eta
 \ll\frac1{Q^2\tau}\ll Q^{-1},
\]
using bounded $f$ and Plancherel. The negative shell is identical. The
bounded range $1\le Q<4$ follows from $|H(a+i\xi)|\ll1/a$ with adjusted
constants. This proves \eqref{eq:high-start}, including off-shell leakage.
\end{proof}

\section{Closure of the weighted estimate}\label{sec:closure}

We now prove Theorem~\ref{thm:energy}. Fix
\begin{equation}\label{eq:threshold}
 0<\sigma\le\sigma_0:=c/16,\qquad
 \eta_0=4\sigma/c,\qquad K=c/(4\sigma).
\end{equation}
Then $\eta_0\le1/4$ and $K\ge4$. Throughout
$\eta_0\le|\xi|\le K$ we have $\sigma\le a(\xi)$: use
$g(\xi)\le3/|\xi|$ for $|\xi|\le1$ and $g(\xi)\le3|\xi|$ for
$|\xi|\ge1$. We claim
\begin{equation}\label{eq:global-Hprime}
 \int_\R|H'(\sigma+i\xi)|^2\dd\xi\ll\sigma^{-1}.
\end{equation}

For $|\xi|\le\eta_0$, differentiating the primitive formula gives
$|H'(s)|\ll\sigma^{-1}+|s|\sigma^{-2}\ll\sigma^{-1}$. The interval has
length $O(\sigma)$, so its squared integral is $O(\sigma^{-1})$.
For $|\xi|\ge K$, \eqref{eq:tail-Hprime} gives the same bound.

In the middle range, $H'=bH-P$ and
\begin{equation}\label{eq:P-L2}
 \int_\R|P(u+i\xi)|^2\dd\xi
 =2\pi\int_0^\infty e^{-2ut}|R(t)|^2\dd t\ll u^{-1}.
\end{equation}
It suffices to prove
\begin{equation}\label{eq:betaH}
 \int_{\eta_0\le|\xi|\le K}|\beta(\sigma,\xi)H(\sigma+i\xi)|^2\dd\xi
 \ll\sigma.
\end{equation}
Apply \eqref{eq:source-bound}, retaining both terms.

The initial-source contribution on a low dyadic shell $[r,2r]$ is, by
\eqref{eq:beta-decay} and \eqref{eq:low-start}, at most
\[
 C\frac{\sigma^2}{r}\left(1+\log\frac r\sigma\right)^2.
\]
Summing upwards from $r=\eta_0$ costs $O(\sigma)$, since
$\sum_{j\ge0}2^{-j}(1+j)^2<\infty$. On a high shell $[Q,2Q]$,
\eqref{eq:high-start} gives at most
\[
 C\sigma^2Q\left(1+\log\frac1{\sigma Q}\right)^2.
\]
Summing downwards from a shell comparable to $K$ again costs $O(\sigma)$.
End shells can be intersected with the middle range. All comparison constants
depend only on the fixed $c$.

For the forcing contribution, abbreviate $a=a(\xi)$ and use
Cauchy--Schwarz in $u$:
\begin{align}
 |\beta|^2\left(\int_\sigma^a|P(u+i\xi)|\dd u\right)^2
 \ll\frac{\sigma^2}{a}\left(1+\log\frac a\sigma\right)^2
                \int_\sigma^a|P(u+i\xi)|^2\dd u.\label{eq:force-CS}
\end{align}
For $\sigma\le u\le a$,
\begin{equation}\label{eq:weight-comparison}
 \frac{(1+\log(a/\sigma))^2}{a}
 \le2\frac{(1+\log(u/\sigma))^2}{u}.
\end{equation}
Indeed $(1+\log x)^2/x$ increases on $[1,e]$ from $1$ to $4/e<2$ and
decreases thereafter. Tonelli's theorem, \eqref{eq:P-L2}, and $a(\xi)\le c$
bound the integral of \eqref{eq:force-CS} over the middle range by
\[
 C\sigma^2\int_\sigma^c\frac{(1+\log(u/\sigma))^2}{u^2}\dd u
 \le C\sigma\int_1^\infty\frac{(1+\log y)^2}{y^2}\dd y\ll\sigma.
\]
This proves \eqref{eq:betaH} and \eqref{eq:global-Hprime}. The forcing was
charged in $L^2$ across frequencies; replacing it everywhere by its pointwise
bound $C/u$ would discard this control.

Finally, Plancherel gives
\[
 \int_0^\infty t^2e^{-2\sigma t}|f(t)|^2\dd t
 =\frac1{2\pi}\int_\R|H'(\sigma+i\xi)|^2\dd\xi\ll\sigma^{-1}.
\]
For $T\ge\sigma_0^{-1}$, set $\sigma=1/T$ and use $e^{-2t/T}\ge e^{-2}$
on $[0,T]$. For $1\le T<\sigma_0^{-1}$, boundedness of $f$ gives
$E(T)\le\norm f_\infty^2T^3/3\ll T$ with the permitted constant.
This completes the abstract proof.

\section{Removing the density hypothesis}\label{sec:wirsing-bridge}

The separate hypothesis \eqref{eq:LD} can be removed if Wirsing's abstract
Tauberian theorem is used as an additional input. The connecting estimate
is global square-integrability of the actual source. We prove it below
from \eqref{eq:B}--\eqref{eq:R}, using only the parts of the preceding
resolvent argument that do not use \eqref{eq:LD}. This gives a second route
to the energy conclusion; the proof in Sections~\ref{sec:scale}--\ref{sec:closure}
continues to give that conclusion without invoking Wirsing.

\begin{lemma}[An interval bound from positivity]\label{lem:weak-density}
Under the hypotheses of Theorem~\ref{thm:energy} other than \eqref{eq:LD},
for every fixed
$A>0$,
\begin{equation}\label{eq:weak-density}
 \nu((t,t+h])\le 2h+\frac{C_A}{t}
 \qquad(t\ge2,\ 0<h\le A).
\end{equation}
In particular, $\nu((t,t+h])\le2h+C/(1+t)$ for $t\ge0$ and $0<h\le1$.
\end{lemma}
\begin{proof}
The function $R$ is locally of bounded variation. Differentiation on the
positive half-line gives the measure identity
\begin{equation}\label{eq:measure-R}
 dR=(t-1)\nu+\nu*\nu+(f-t-B-R)\dd t.
\end{equation}
Indeed $d(f*\nu)=df*\nu-\nu$, with the last term supplied by $f(0)=-1$.
Substitute $df=\nu-(1+f)\dd t$ and $f*\nu=R-tf$ to obtain
\eqref{eq:measure-R}. Positivity, integrated over $(t,t+h]$, now gives
\[
 (t-1)\nu((t,t+h])
 \le R(t+h)-R(t)+\int_t^{t+h}(u+B(u)+R(u)-f(u))\dd u
 \le 2th+h^2+Ch+2\norm R_\infty.
\]
Division by $t-1$ proves \eqref{eq:weak-density}; the bounded initial
interval is absorbed by increasing the constant.
\end{proof}

\begin{theorem}[Square-integrability without a separate density input]
\label{thm:weak-L2}
Assume all the base hypotheses of Theorem~\ref{thm:energy}, including
\eqref{eq:B}--\eqref{eq:R}, but omit \eqref{eq:LD}. Then
\begin{equation}\label{eq:unweighted-L2}
 \int_0^\infty |f(t)|^2\dd t<\infty.
\end{equation}
The integral is bounded in terms of the base hypothesis bounds.
\end{theorem}
\begin{proof}
Write $\ell_\sigma=\log(e/\sigma)$ for $0<\sigma\le1$.
Use Lemma~\ref{lem:weak-density} in the proof of
Lemma~\ref{lem:translations}, with $k=0$ and $q(t)=C/(1+t)$.
The two bounds for the source increment are now
$Ch+\nu((t,t+h])$ and $Ch+C/(1+t)$. Multiplication and Tonelli give
\begin{equation}\label{eq:weak-translation}
 \norm{U_{\sigma,0}(\cdot-h)-U_{\sigma,0}}_2^2
 \ll h^2/\sigma+h\ell_\sigma\qquad(0<h\le1).
\end{equation}
For the two terms containing $q$, use
\[
 \int_0^\infty\frac{e^{-2\sigma u}}{1+u}
       \bigl(\dd u+\nu(\mathrm du)\bigr)\ll\ell_\sigma.
\]
The measure estimate follows by dividing into unit intervals and using
$B(u)=u+O(1)$. Weight changes and the strip at zero cost
$O(h^2\sigma+h)$, within \eqref{eq:weak-translation}.
Averaging the Fourier translation identity as before yields
\begin{equation}\label{eq:weak-tail}
 \int_{|\xi|\ge K}|H(\sigma+i\xi)|^2\dd\xi
 \ll\frac1{\sigma K^2}+\frac{\ell_\sigma}{K},\qquad K\ge2.
\end{equation}

The scale equation, two-frequency exclusion, and source resolvent in
Sections~\ref{sec:scale}--\ref{sec:resolvent} used only the base hypotheses.
Retain their $c$, $a(\xi)$, and $H_*(\xi)$. The low-shell bound
\eqref{eq:low-start} also used only boundedness of $F$. On a high shell,
repeat the Poisson argument of Lemma~\ref{lem:start}, now with
\eqref{eq:weak-tail} on the line $\tau=c/(8Q)$. The enlarged shell costs
$O(\log(2Q)/Q)$, and the off-shell leakage is still $O(1/Q)$.
The bounded ratio of widths in \eqref{eq:poisson-ratio} is unchanged. Thus
\begin{equation}\label{eq:weak-start}
 \int_{Q\le|\xi|\le2Q}|H_*(\xi)|^2\dd\xi
 \ll\frac{\log(2Q)}Q\qquad(Q\ge1).
\end{equation}
The tail calculation uses $Q\ge4$; the remaining compact range
$1\le Q<4$ follows from boundedness of $f$ and a positive lower bound
for $a(\xi)$ there.
Both this high-shell series and the low-shell series $O(r)$ converge
dyadically, so $\norm{H_*}_2\ll1$.

Fix $0<\sigma\le c/16$ and let
$D_\sigma=\{4\sigma/c\le|\xi|\le c/(4\sigma)\}$.
Here $\sigma\le a(\xi)$, so \eqref{eq:source-bound} applies. Its forcing
term is bounded in $L^2$ by Minkowski's inequality:
\begin{align}
 \left\|\ind_{D_\sigma}\int_\sigma^{a(\xi)}|P(u+i\xi)|\dd u\right\|_2
 &\le\int_\sigma^c
       \norm{\ind_{D_\sigma}\ind_{\{u\le a(\xi)\}}P(u+i\xi)}_2\dd u
 \notag\\
 &\ll\int_\sigma^c u^{-1/2}\dd u\ll1.\label{eq:weak-force}
\end{align}
The estimate on $P$ is precisely \eqref{eq:P-L2}, which requires only
bounded $R$. Hence $\int_{D_\sigma}|H(\sigma+i\xi)|^2\dd\xi\ll1$.
On $|\xi|<4\sigma/c$, the bounded primitive gives
$|H(s)|\ll|s|/\sigma\ll1$, at squared-integral cost $O(\sigma)$.
On $|\xi|>c/(4\sigma)$, \eqref{eq:weak-tail} costs
$O(\sigma\ell_\sigma)$. We have proved
\[
 \sup_{0<\sigma\le c/16}\int_\R|H(\sigma+i\xi)|^2\dd\xi<\infty.
\]
Plancherel and monotone convergence as $\sigma\downarrow0$ prove
\eqref{eq:unweighted-L2}. This proof has not used the weighted energy
theorem or any pointwise PNT consequence.
\end{proof}

\begin{corollary}[A bootstrap using Wirsing]\label{cor:wirsing-bootstrap}
Under the hypotheses of Theorem~\ref{thm:weak-L2}, Wirsing's Satz~2
\cite[Section~5, p.~15]{wirsing2} gives
\begin{equation}\label{eq:wirsing-bootstrap}
 f(t)=O(t^{-1}),\qquad E(T)=O(T).
\end{equation}
In the arithmetic instance, this yields $\psi(x)=x+O(x/\log x)$.
This corollary depends on Wirsing's Satz~2, which he stated without a
detailed proof \cite[p.~15]{wirsing2}. The direct proof of the weaker
pointwise rate in Corollary~\ref{cor:PNT} does not depend on it.
\end{corollary}
\begin{proof}
First $f(t)\to0$. The source equation and Lemma~\ref{lem:weak-density}
give $|f(t+h)-f(t)|\le Ch+C/(1+t)$ for $0<h\le1$.
If $|f(t_j)|\ge\varepsilon$ at arbitrarily large $t_j$, this inequality
gives disjoint forward intervals of a fixed positive length on which
$|f|\ge\varepsilon/2$, contradicting \eqref{eq:unweighted-L2}.
This also treats the right-continuous values at atoms.

Put $g=-f$ and fix $\delta>0$ with $\gamma:=1+\delta<\sqrt2$.
For all sufficiently large $t$, $|f(u)|\le\delta$ for $u\ge t$.
For each fixed $A$ the exact increment identity and
\eqref{eq:weak-density} give
\[
 -(1+\delta)h\le f(t+h)-f(t)
 =\nu((t,t+h])-h-\int_t^{t+h}f(u)\dd u
 \le(1+\delta)h+C_A/t\qquad(0<h\le A).
\]
Taking $A=\max(1,2\norm f_\infty)$ handles all larger $h$ by boundedness.
This holds for $t\ge t_\delta$, where $t_\delta\ge2$ depends on the pair.
For $t<t_\delta$, $|g(t+h)-g(t)|\le2\norm f_\infty\le2\norm f_\infty t_\delta/t$.
Consequently
\begin{equation}\label{eq:bridge-fluctuation}
 |g(t+h)-g(t)|\le\gamma h+O(t^{-1}),\qquad h>0.
\end{equation}
Also $\int_{(0,T]}|dg|\le B(T)+\int_0^T|1+f|\dd t=O(T)$.
Finally, expanding $f*\nu$ using the source equation gives the exact identity
\begin{equation}\label{eq:bridge-selfconv}
 tg-g*dg=F+f*f-R=O(1),
\end{equation}
because $|f*f|\le\norm f_2^2$ and both $F$ and $R$ are bounded.
Thus $g$ has the total-variation bound, fluctuation bound with
$\gamma<\sqrt2$, and self-convolution equation required by Satz~2,
with Wirsing's exponent $\beta=1$ (unrelated to the variable $\beta$ of
Section~\ref{sec:scale}).

Wirsing uses a zero extension at $0-$ and allows the countably many
points where a Stieltjes convolution is undefined. Adding the initial
atom to our derivative convention changes $g*dg$ by $g(0)g(t)$, a bounded
term. If $J$ is the countable jump set of $g$, simultaneous interior
jumps occur only for $t\in J+J$. Outside this countable set the
Riemann--Stieltjes and measure convolutions agree; the printed exception
covers the remaining points. The upper-endpoint value agrees because
$g$ is right-continuous at zero. Therefore his theorem applies
and proves \eqref{eq:wirsing-bootstrap}. The pointwise implied constant
here is asserted for each fixed admissible pair; uniform dependence solely
on the original hypothesis bounds is not claimed by this bootstrap.
\end{proof}

\section{Classical arithmetic inputs and the coupled source}\label{sec:inputs}

We record the classical arithmetic ingredients to make clear where the
analytic hypotheses come from. None is a PNT remainder assumption.

\begin{lemma}[Chebyshev and the harmonic prime-power sum]\label{lem:elementary}
For $x\ge2$,
\begin{equation}\label{eq:elementary}
 \vartheta(x)\ll x,\qquad \psi(x)\ll x,\qquad
 \sum_{n\le x}\frac{\Lambda(n)}n=\log x+O(1).
\end{equation}
Also $\pi(x)\ll x/\log x$.
\end{lemma}
\begin{proof}
Every prime in $(n,2n]$ divides $\binom{2n}{n}$, so
$\vartheta(2n)-\vartheta(n)\le2n\log2$. Summing at powers of two and using
monotonicity gives $\vartheta(x)\ll x$. The identity
$\psi(x)=\sum_{a\ge1}\vartheta(x^{1/a})$ then gives $\psi(x)\ll x$:
the $a\ge2$ terms are $O(\sqrt{x}\log x)=O(x)$.

From $\sum_{d\mid n}\Lambda(d)=\log n$,
\[
 \sum_{d\le x}\Lambda(d)\left\lfloor\frac xd\right\rfloor
 =\sum_{n\le x}\log n=x\log x-x+O(\log(2x)).
\]
Replacing the floor by $x/d$ costs at most $\psi(x)=O(x)$, which proves
the harmonic estimate. Finally primes greater than $\sqrt{x}$ each
contribute at least $(\log x)/2$ to $\vartheta(x)$; the remaining count
is at most $\sqrt{x}$. This proves the bound on $\pi$.
\end{proof}

\begin{lemma}[Selberg's symmetry formula]\label{lem:selberg}
For $x\ge2$,
\begin{equation}\label{eq:selberg}
 \sum_{n\le x}\bigl(\Lambda(n)\log n+(\Lambda*_D\Lambda)(n)\bigr)
 =2x\log x+O(x).
\end{equation}
\end{lemma}
\begin{proof}
Here is the elementary inversion argument, in the form discussed
in \cite[Section 1]{tao}. Put
$\Lambda_2(n)=\sum_{d\mid n}\mu(d)\log^2(n/d)$.
The identity $\Lambda_2=\Lambda\log+ \Lambda*_D\Lambda$ follows by
applying logarithmic differentiation to $\mathbf1*_D\Lambda=\log$
and then convolving with $\mu$.

Integral comparison and partial summation give constants $c_1,c_2,c_3$ with
\begin{align*}
 \frac1y\sum_{n\le y}\log^2n
   &=\log^2y-2\log y+2+O\!\left(\frac{\log^2(2+y)}y\right),\\
 \sum_{n\le y}\frac1n&=\log y+c_1+O(y^{-1}),\\
 \sum_{n\le y}\frac{\log(y/n)}n
   &=\tfrac12\log^2y+c_2\log y+c_3
       +O\!\left(\frac{\log(2+y)}y\right).
\end{align*}
For the last line one may integrate the preceding harmonic estimate;
its $O(1/y)$ remainder has a convergent integral against $dy/y$.
Consequently there are constants $c_4,c_5$ such that
\[
 \frac1y\sum_{n\le y}\log^2n
 =\sum_{m\le y}\frac{2\log(y/m)+c_4+c_5/(y/m)}m
   +O\!\left(\frac{\log^2(2+y)}y\right).
\]
Insert $y=x/d$ and sum with coefficient $\mu(d)/d$.
The error is bounded by
$x^{-1}\sum_{d\le x}\log^2(2+x/d)=O(1)$, as follows by dyadic grouping.
Grouping the main term by $k=dm$ and using
$\sum_{d\mid k}\mu(d)=\ind_{k=1}$ leaves
$2\log x+c_4+c_5/x$. Multiplication by $x$ proves the lemma.
\end{proof}

\begin{proposition}\label{prop:arithmetic-source}
The source and measure in \eqref{eq:arith-source} satisfy
\eqref{eq:B}--\eqref{eq:R}.
\end{proposition}
\begin{proof}
The jump of $f$ at $\log n$ is $\Lambda(n)/n$, and between jumps
$f'=-(1+f)$. Also $f(0)=-1$ and $\nu$ has no atom at zero. Thus
\eqref{eq:source} holds with the stated endpoint convention.
Lemma~\ref{lem:elementary} supplies bounded $f$ and $B(t)=t+O(1)$.

Put $x=e^t$. Partial summation and Chebyshev give
\[
 \sum_{n\le x}\Lambda(n)\log n
   =\psi(x)\log x-\int_1^x\frac{\psi(u)}u\dd u
   =\psi(x)\log x+O(x).
\]
Moreover
\[
 \frac1x\sum_{n\le x}(\Lambda*_D\Lambda)(n)
 =\sum_{k\le x}\frac{\Lambda(k)}k
       \bigl(1+f(\log x-\log k)\bigr)
 =B(t)+(f*\nu)(t).
\]
Divide \eqref{eq:selberg} by $x$ and substitute these identities.
The two main terms $t$ cancel against $2t$, leaving
$tf(t)+(f*\nu)(t)=O(1)$ uniformly for $t\ge\log2$.
The bounded initial interval is harmless.
\end{proof}

\section{The local-density hypothesis from a quadratic sieve}\label{sec:sieve}

The sieve below is Selberg's quadratic sieve \cite{selberg-sieve}, in a
coarse form. Sharp constants in the resulting Brun--Titchmarsh inequality
are due to Montgomery and Vaughan \cite{mv}; we need only the order of
magnitude.

\begin{lemma}[A coarse uniform interval sieve]\label{lem:sieve}
For $X\ge1$ and $Y\ge2$,
\begin{equation}\label{eq:interval-sieve}
 \pi(X+Y)-\pi(X)\ll\frac{Y}{\log Y}.
\end{equation}
\end{lemma}
\begin{proof}
For an integer $z\ge1$, write
$G_z=\sum_{\ell\le z}\mu(\ell)^2/\varphi(\ell)$.
Choose real weights $\lambda_d$ supported on $d\le z$, with $\lambda_1=1$.
The identity $\gcd(d,e)=\sum_{\ell\mid d,\,\ell\mid e}\varphi(\ell)$ gives
\begin{equation}\label{eq:sieve-form}
 Q(\lambda):=\sum_{d,e\le z}\frac{\lambda_d\lambda_e}{\lcm(d,e)}
 =\sum_{\ell\le z}\varphi(\ell)
       \left(\sum_{\substack{d\le z\\\ell\mid d}}\frac{\lambda_d}{d}\right)^2.
\end{equation}
Put $x_\ell=\sum_{\ell\mid d,\,d\le z}\lambda_d/d$. Inversion gives
$\lambda_1=\sum_{\ell\le z}\mu(\ell)x_\ell$.
Cauchy--Schwarz therefore proves $Q(\lambda)\ge1/G_z$, with equality at
$x_\ell=\mu(\ell)/(\varphi(\ell)G_z)$. Inverting these minimizing values,
\[
 \lambda_d=\frac{\mu(d)d}{\varphi(d)G_z}
    \sum_{\substack{k\le z/d\\(k,d)=1}}\frac{\mu(k)^2}{\varphi(k)}.
\]
Since the inner sum is at most $G_z$, $|\lambda_d|\le d/\varphi(d)\le d$ and
$\sum_{d\le z}|\lambda_d|\le z^2$.

For each prime $n>z$, $\sum_{d\mid n}\lambda_d=1$.
Summing the nonnegative square of this divisor sum on $(X,X+Y]$ gives
\begin{equation}\label{eq:sieve-count}
 \pi(X+Y)-\pi(X)\le z+\frac{Y}{G_z}+O(z^4).
\end{equation}
The error comes from the uniform $O(1)$ error in counting multiples of
$\lcm(d,e)$.

At least one quarter of the integers through any integer $n$ are
squarefree: the nonsquarefree count is at most
$n\sum_{k\ge2}k^{-2}<3n/4$.
Partial summation gives
$G_z\ge\sum_{\ell\le z}\mu(\ell)^2/\ell\gg\log z$ for large $z$.
Take $z=\lfloor Y^{1/8}\rfloor$. The $z$ and $z^4$ terms are absorbed by
$CY/\log Y$ for large $Y$, and the bounded remaining range is absorbed
by increasing $C$. This proves \eqref{eq:interval-sieve}.
\end{proof}

\begin{proposition}\label{prop:LD-arithmetic}
The prime-power measure in \eqref{eq:arith-source} satisfies
\eqref{eq:LD}.
\end{proposition}
\begin{proof}
For large $t$ and $e^{-t/2}\le h\le1$, put
$X=e^t$, $Y=e^t(e^h-1)$. Then $\log Y\ge t/2$ and $Y\ll he^t$.
Lemma~\ref{lem:sieve} gives
\[
 \sum_{e^t<p\le e^{t+h}}\frac{\log p}{p}\ll h.
\]
Prime powers of exponent at least two contribute
$O((1+t)^2e^{-t/2})$: their count through $e^{t+1}$ is
$O((1+t)e^{(t+1)/2})$, and each weight in the interval is at most
$(t+1)e^{-t}$.

For $0<h<e^{-t/2}$, enlarge the interval to logarithmic width $e^{-t/2}$.
Its prime contribution is $O(e^{-t/2})$, absorbed by the same error
majorant. A larger constant covers the bounded initial range of $t$.
Thus one may take $q(t)=C(1+t)^2e^{-t/2}$.
\end{proof}

Propositions~\ref{prop:arithmetic-source} and~\ref{prop:LD-arithmetic}
discharge every arithmetic hypothesis of Theorem~\ref{thm:energy}.
The sieve is an upper-density estimate; it assumes no asymptotic prime
count in short intervals.

\section{The exact sparse Mertens certificate}\label{sec:certificate}

Let $N\ge12$ be an integer. Write $\mathcal S_N$ for the squarefree
positive integers at most $N$, and
\[
 \mathcal T_N=\{n\in\mathcal S_N:n\text{ odd},\ N/2<n\le N\},\qquad
 \mathcal C_N=\{n\in\mathcal T_N:n\text{ composite}\}.
\]
The admissible labels are
\[
 \mathcal B_N=\{b:5\le b\le N/3,\ b\text{ squarefree},\ (b,6)=1\}.
\]
For any real coefficient vector $\alpha=(\alpha_b)_{b\in\mathcal B_N}$, set
\begin{align}
 (A\alpha)(n)&=\sum_{\substack{b\in\mathcal B_N\\b\mid n}}\alpha_b,
        \qquad n\in\mathcal T_N,\label{eq:A}\\
 (D\alpha)(m)&=\sum_{\substack{b\in\mathcal B_N\\N/2<bm\le N\\(b,m)=1}}
                       \mu(b)\alpha_b,
        \qquad m\text{ odd and squarefree}.\label{eq:D}
\end{align}
The second vector vanishes for $m>N/5$. The matrices are sparse incidence
matrices of the same multiplication pairs, with the label sign retained
in the cofactor block.

\begin{proposition}[Exact finite identity]\label{prop:certificate}
For every such $\alpha$,
\begin{equation}\label{eq:finite-identity}
 M(N)=\sum_{n\in\mathcal T_N}\mu(n)(1-(A\alpha)(n))
       +\sum_{\substack{m\le N/5\\m\ \mathrm{odd,\ squarefree}}}
                   \mu(m)(D\alpha)(m).
\end{equation}
Consequently
\begin{equation}\label{eq:certificate}
 |M(N)|\le \mathfrak B_N(\alpha):=
   \pi(N)-\pi(N/2)+\norm{1-A\alpha}_{1,\mathcal C_N}
                      +\norm{D\alpha}_1.
\end{equation}
\end{proposition}
\begin{proof}
Pair $n=m$ and $n=2m$ for odd squarefree $m$ with $2m\le N$.
Their M\"obius values cancel, so
$M(N)=\sum_{n\in\mathcal T_N}\mu(n)$.
Expanding $\sum_{n\in\mathcal T_N}\mu(n)(A\alpha)(n)$ and writing $n=bm$
gives the second sum in \eqref{eq:finite-identity}, since
$\mu(bm)=\mu(b)\mu(m)$ on its support.
Take absolute values. No admissible label divides a prime in $(N/2,N]$,
so each such prime contributes exactly one to the first residual norm.
This yields \eqref{eq:certificate}.
\end{proof}

The target estimate is
\begin{equation}\label{eq:target-six}
 \norm{1-A\alpha}_{1,\mathcal C_N}+\norm{D\alpha}_1
       \ll \frac{N}{\log N}.
\end{equation}
No normalization or parity-balance constraint on $\alpha$ is required
by the exact identity.

\section{Explicit weights and the cofactor reduction}\label{sec:weights}

Put $L=\log N$. Retaining only admissible labels, choose
\begin{equation}\label{eq:weights}
 \alpha_p=\frac{(\log p)^2}{L^2},\qquad
 \alpha_{pq}=\frac{2\log p\log q}{L^2}\quad(5\le p<q),\qquad
 \alpha_b=0\ \text{otherwise}.
\end{equation}
These are the squarefree prime and semiprime terms in the classical second
generalized von Mangoldt weights.

\begin{lemma}[Coverage and coefficient mass]\label{lem:coverage}
For the coefficients \eqref{eq:weights},
\begin{equation}\label{eq:coverage}
 \norm{1-A\alpha}_{1,\mathcal C_N}\ll N/L.
\end{equation}
If $\rho_b=\#\{n\in\mathcal S_N:b\mid n\}/|\mathcal S_N|$, then
$\sum_b\rho_b|\alpha_b|\le1$.
\end{lemma}
\begin{proof}
For $n\in\mathcal C_N$, let $\epsilon_n=\ind_{3\mid n}$. Expanding the square
of the sum of logarithms of its prime factors at least five gives exactly
\begin{equation}\label{eq:coverage-exact}
 (A\alpha)(n)=\frac{(\log n-\epsilon_n\log3)^2}{L^2}
 -\ind_{\substack{n=pq\\3\nmid n}}\frac{2\log p\log q}{L^2}.
\end{equation}
Every relevant prime or two-prime divisor is at most $N/3$ except the
whole number $n=pq$ when $3\nmid n$; this is the subtracted term.
The first term has deficit at most $2\log6/L$, since
$N/2<n\le N$. Also, by Chebyshev and partial summation,
\[
 \sum_{\substack{p<q\\pq\le N}}2\log p\log q
 \le2\sum_{p\le\sqrt N}\log p\,\vartheta(N/p)
 \ll N\sum_{p\le\sqrt N}\frac{\log p}{p}\ll NL.
\]
This proves \eqref{eq:coverage}. For every squarefree $n\le N$, even or odd,
the full retained divisor sum is nonnegative and at most
$(\sum_{p\mid n,\ p\ge5}\log p)^2/L^2\le1$.
Averaging that inequality over $\mathcal S_N$ proves the mass bound.
\end{proof}

For an integer $Q\ge1$ and real $x\ge1$, define
\[
 G_Q(x)=-\sum_{\substack{p\le x\\(p,Q)=1}}(\log p)^2
     +2\sum_{\substack{p<q,\ pq\le x\\(pq,Q)=1}}\log p\log q,
 \qquad G=G_1.
\]
For odd squarefree $m\ge3$, the cutoff $b\le N/3$ is automatic when
$b\le N/m$. Hence
\begin{equation}\label{eq:flux-G}
 (D\alpha)(m)=\frac{G_{6m}(N/m)-G_{6m}(N/(2m))}{L^2}.
\end{equation}
The row $m=1$ is zero because all labels are at most $N/3$.

\begin{lemma}[Coprimality removal]\label{lem:coprime}
The aggregate error in replacing both occurrences of $G_{6m}$ by $G$
in \eqref{eq:flux-G}, before division by $L^2$, is $O(NL)$.
\end{lemma}
\begin{proof}
The union bound and Chebyshev give
\[
 |G(x)-G_Q(x)|
 \le\sum_{\substack{p\mid Q\\p\le x}}(\log p)^2
      +2\sum_{p\mid Q}\log p\,\vartheta(x/p).
\]
Enlarge the sum over $m$ to all integers $3\le m\le N/5$.
For $p=2,3$ the total is $O(NL)$. The terms with $p\mid m$ in the second
sum are bounded by
\[
 CN\sum_{m\le N/5}\frac1m\sum_{p\mid m}\frac{\log p}{p}
 \ll NL\sum_p\frac{\log p}{p^2}\ll NL.
\]
For the first sum, $p\mid m$ and $p\le N/m$ imply $p^2\le N$, and
there are at most $N/p^2$ eligible $m$. Their total is
$O(N\sum_p(\log p)^2/p^2)=O(N)$. Both series converge by comparison
with sums over the integers. The same estimates apply at $N/(2m)$.
\end{proof}

\begin{lemma}[The signed prime/semiprime sum]\label{lem:G}
For $x\ge2$,
\begin{equation}\label{eq:G-source}
 G(x)=-2x\log x\,f(\log x)+O(x).
\end{equation}
\end{lemma}
\begin{proof}
Replacing prime sums by prime-power sums gives
\[
 G(x)=\sum_{n\le x}\bigl((\Lambda*_D\Lambda)(n)-\Lambda(n)\log n\bigr)+O(x).
\]
For completeness, the convolution terms with at least one proper prime
power cost at most
$2\sum_{p^a\le x,\ a\ge2}\log p\,\psi(x/p^a)
 \ll x\sum_{p,a\ge2}\log p/p^a=O(x)$.
The prime diagonal and the proper-power terms in $\Lambda(n)\log n$
also cost $O(x)$; the crude bound $O(\sqrt x\,\log^3(2x))=O(x)$
suffices. Selberg's formula and partial summation therefore give
\[
 G(x)=2x\log x-2\psi(x)\log x+O(x),
\]
which is \eqref{eq:G-source}.
\end{proof}

\section{Sampling, the tail estimate, and the arithmetic conclusions}
\label{sec:transfer}

\begin{lemma}[Sampling with a logarithmic weight]\label{lem:sampling}
For real $X\ge4$, put $L_X=\log X$, $J_X=\lfloor X/2\rfloor$, and
\[
 \Phi_X=\sum_{m\le J_X}\frac{\log(X/m)}m\,|f(\log(X/m))|.
\]
Then
\begin{equation}\label{eq:sampling}
 \Phi_X\le\int_0^{L_X}t|f(t)|\dd t+O(L_X).
\end{equation}
\end{lemma}
\begin{proof}
Set $g_X(u)=\psi(X/u)/X-1/u=f(\log(X/u))/u$.
For $m\le u\le m+1$,
\[
 |g_X(m)-g_X(u)|
 \le\frac{\psi(X/m)-\psi(X/(m+1))}{X}+\frac1{m(m+1)}.
\]
Multiply by $\log(X/m)$ and integrate over this unit interval.
The first errors telescope at cost at most $L_X\psi(X)/X=O(L_X)$;
the second errors also cost $O(L_X)$.
Replacing $\log(X/m)$ under the integral by $\log(X/u)$ costs $O(1)$,
since $|g_X(u)|\ll1/u$ and $\log(u/m)\le1/m$.
The substitution $t=\log(X/u)$ gives an integral over a subinterval of
$[0,L_X]$ because $J_X+1\le X$. Enlarging that interval proves the claim.
\end{proof}

The arithmetic instance of Theorem~\ref{thm:energy} and Cauchy--Schwarz give
\begin{equation}\label{eq:Phi-bound}
 \Phi_X\le\sqrt{L_X E(L_X)}+O(L_X)\ll L_X.
\end{equation}
In particular there is no additional $\log\log X$ loss in this transfer.

\begin{theorem}[Closure of the sparse certificate]\label{thm:certificate}
The explicit coefficients \eqref{eq:weights} satisfy \eqref{eq:target-six}
uniformly for integers $N\ge12$. Consequently $|M(N)|\ll N/\log N$.
\end{theorem}
\begin{proof}
Lemma~\ref{lem:G} gives, for $3\le m\le N/5$,
\begin{align*}
 |G(N/m)-G(N/(2m))|
 \le\frac Nm\bigl(&2\log(N/m)|f(\log(N/m))|\\
                 &+\log(N/(2m))|f(\log(N/(2m)))|+C\bigr).
\end{align*}
Summing and using \eqref{eq:Phi-bound}, once with $X=N$ and once with
$X=N/2$, yields
\begin{equation}\label{eq:cofactor-tail}
 \sum_{\substack{3\le m\le N/5\\m\ \mathrm{odd,\ squarefree}}}
       |G(N/m)-G(N/(2m))|\ll N\log N.
\end{equation}
The $O(N/m)$ errors sum to $O(N\log N)$ as well.
Now \eqref{eq:flux-G} and Lemma~\ref{lem:coprime}, after division by $L^2$,
give $\norm{D\alpha}_1\ll N/L$. Lemma~\ref{lem:coverage} gives the other
term in \eqref{eq:target-six}.
Finally apply \eqref{eq:certificate} and $\pi(N)\ll N/\log N$.
\end{proof}

\begin{corollary}[A direct quantitative PNT]\label{cor:PNT}
For $x\ge2$,
$\psi(x)=x+O(x/(\log x)^{1/3})$.
\end{corollary}
\begin{proof}
The source equation and \eqref{eq:LD} imply
$|f(t+h)-f(t)|\le C_0h+q(t)$ for $0<h\le1$.
Increase $C_0$ so that $|f(t)|/(4C_0)\le1$.
Write $A=|f(t)|$. If $A\ge4q(t)$, then $|f(t+h)|\ge A/2$ for
$0\le h\le h_0:=A/(4C_0)$. Hence, for $t\ge1$,
\[
 E(t+1)\ge\int_t^{t+h_0}u^2|f(u)|^2\dd u
           \ge\frac{t^2A^3}{16C_0}.
\]
The energy bound gives $A\ll t^{-1/3}$. If $A<4q(t)$, the same conclusion
follows from the exponential majorant of $q$.
The range $\log2\le t<1$ follows from boundedness of $f$.
Putting $t=\log x$ proves the corollary, including jump locations under
the right-continuous convention.
\end{proof}

\begin{remark}[The certificate's floor]\label{rem:floor}
For every admissible coefficient vector,
$\mathfrak B_N(\alpha)\ge\pi(N)-\pi(N/2)$.
The PNT just proved, followed by partial summation, implies
$\pi(N)-\pi(N/2)\sim N/(2\log N)$. Thus the absolute-value certificate
itself cannot give $o(N/\log N)$, even with better control of its other
terms. This is a limitation of the certificate, not a lower bound for
$|M(N)|$. A stronger rate requires retaining further signed cancellation
or changing the final certificate.
\end{remark}

\section{Finite eigenstructure, singular vectors, and scope}\label{sec:finite}

This section records the finite linear-algebra formulation. Its identities
are exact, but the analytic theorem above should not be read as proving
every natural finite inverse estimate.

\subsection{Similarity to a diagonal operator}
Allow a real endpoint $N\ge2$ and put $J=\lfloor N/2\rfloor$,
$t_m=\log(N/m)$. On vectors indexed by $1\le m\le J$, define
\begin{align}
 (S_Nr)_m&=t_mr_m+
       \sum_{\substack{k\ge2\\mk\le J}}\frac{\Lambda(k)}k\,r_{mk},
       \label{eq:finite-S}\\
 (Z_Nv)_m&=\sum_{mk\le J}\frac{v_{mk}}k,\qquad
 (Z_N^{-1}v)_m=\sum_{mk\le J}\frac{\mu(k)}k\,v_{mk},\qquad
 Q_N=\diag(t_m).\label{eq:finite-Z}
\end{align}
Finite divisor inversion gives the inverse in \eqref{eq:finite-Z}, and
\begin{equation}\label{eq:similarity}
 S_N=Z_N^{-1}Q_NZ_N.
\end{equation}
For instance the coefficient of $v_{mk}/k$ in
$Q_NZ_N-Z_NQ_N$ is $\log k$; the identity
$\sum_{d\mid k}\Lambda(d)=\log k$ proves $Z_NS_N=Q_NZ_N$.

Thus the eigenvalues are the distinct positive numbers $t_m\ge\log2$.
Right eigenvectors are $Z_N^{-1}e_m$ and left eigenvectors are
$e_m^{\mathsf T}Z_N$, with their natural biorthogonal normalization.
This exact eigenstructure does not control the inverse norm in a
scale-independent way: the similarity is not asserted to be uniformly
well-conditioned.

For the actual vector
$r_m=f(\log(N/m))=\psi(N/m)/(N/m)-1$, the elementary arithmetic inputs give
\begin{equation}\label{eq:finite-forcing}
 \norm{S_Nr}_\infty\ll1.
\end{equation}
This follows from \eqref{eq:R} by omitting $k>N/(2m)$, where
$f(\log(N/(mk)))=-1$. The omitted harmonic prime-power sum is $O(1)$
on a multiplicative interval of ratio two. The actual vector also satisfies
\begin{equation}\label{eq:counting-constraints}
 0\le1+r_m\le C_1,\qquad
 \frac{1+r_m}{m}\ge\frac{1+r_{m+1}}{m+1}\quad(m<J),
\end{equation}
because $(1+r_m)/m=\psi(N/m)/N$.

\subsection{The exact SVD expression and an obstruction}
Let $\mathsf W=\diag(1/m)$ and
\[
 \mathsf T_N=\mathsf W^{1/2}S_NQ_N^{-1}\mathsf W^{-1/2}.
\]
If $\mathsf T_N=U\Sigma V^*$ has singular values $s_j>0$ and left
singular vectors $u_j$, and $y=\mathsf W^{1/2}S_Nr$, then
\begin{equation}\label{eq:svd}
 \sum_{m\le J}\frac{t_m^2|r_m|^2}{m}
   =\sum_{j=1}^J\frac{|\langle u_j,y\rangle|^2}{s_j^2}.
\end{equation}
Indeed $\mathsf T_N(\mathsf W^{1/2}Q_Nr)=y$.
Equation~\eqref{eq:svd} identifies the required interaction between small
singular values and the distinguished source. It supplies no lower bound
on every singular value.

There is a concrete obstruction to an unrestricted source estimate.
Let $\lambda(n)=(-1)^{\Omega(n)}$ be the Liouville function and take
$r_m=\lambda(m)$. Complete multiplicativity and
Lemma~\ref{lem:elementary} give, for $y\ge1$,
\[
 \sum_{k\le y}\frac{\Lambda(k)\lambda(k)}k
 =-\sum_{k\le y}\frac{\Lambda(k)}k
       +2\sum_{p^{2a}\le y}\frac{\log p}{p^{2a}}
 =-\log y+O(1).
\]
Consequently $\norm{S_Nr}_\infty=O(1)$, whereas
\begin{align}
 \sum_{m\le J}\frac{t_m|r_m|}{m}
 &=\tfrac12(\log N)^2+O(\log N),\label{eq:liouville-obstruction}\\
 \sum_{m\le J}\frac{t_m^2|r_m|^2}{m}
 &=\tfrac13(\log N)^3+O((\log N)^2).\notag
\end{align}
Integral comparison proves the two sums.
Thus bounded forcing alone cannot yield the desired source bounds.
The Liouville vector meets the amplitude constraint but, for large $N$,
violates the monotonicity in \eqref{eq:counting-constraints}.

A sufficient spectral-tail estimate would be, for some fixed $\delta>0$,
\begin{equation}\label{eq:svd-tail-open}
 \sum_{s_j\le s}|\langle u_j,y\rangle|^2
 \le C(\log N)s^{2+\delta},\qquad0<s\le1.
\end{equation}
Together with $\norm y_2^2=O(\log N)$, dyadic summation in
\eqref{eq:svd} would give a discrete weighted energy of order $\log N$.
The exponent $2$ without a further summable factor does not suffice for
that inference. Estimate~\eqref{eq:svd-tail-open} is a sufficient target,
not a theorem established here.

\subsection{The stronger constrained inverse remains open}
The stronger statement would be that every real vector satisfying
\eqref{eq:counting-constraints} and
$\norm{S_Nr}_\infty\le C_0$ obeys
\begin{equation}\label{eq:constrained-open}
 \sum_{m\le J}\frac{t_m|r_m|}{m}\le C(C_0,C_1)\log N.
\end{equation}
Theorem~\ref{thm:energy} and the sampling argument prove the required bound
for the actual arithmetic source. They do not prove
\eqref{eq:constrained-open} for all artificially constructed vectors:
the analytic proof also uses the common positive measure and
\eqref{eq:LD}.

There is a finite, directly verifiable dual sufficient condition for
\eqref{eq:constrained-open}. Define
$(\mathcal K r)_m=r_m/m-r_{m+1}/(m+1)$ and $d_m=1/[m(m+1)]$ for $m<J$.
For every sign vector $\epsilon\in\{-1,1\}^J$ put $c_m=\epsilon_mt_m/m$.
It would suffice to construct $a\in\R^J$ and nonnegative vectors $u,v,z$
with $z\in\R^{J-1}$ such that
\begin{align}
 c&=S_N^{\mathsf T}a+u-v-\mathcal K^{\mathsf T}z,\label{eq:dual-open}\\
 C_0\norm a_1+(C_1-1)\sum_m u_m+\sum_m v_m+\sum_m z_md_m
    &\le C\log N.\notag
\end{align}
Pair the equality with $r$, use $\mathcal K r\ge-d$, $r\le C_1-1$, and
$-r\le1$, and then maximize over $\epsilon$.
This implication is just a finite inequality; no separation theorem is
needed. A uniform construction of these multipliers has not been supplied.

\section{Exact recursions and what they do not yet give}\label{sec:recursion}

\subsection{Restriction to divisibility rays}
For an integer $d\ge2$, put $X=N/d$ and $r^{(d)}_k=r_{dk}$ for
$k\le\lfloor X/2\rfloor$ (the identity below is trivial when $X<4$). From \eqref{eq:finite-S},
\begin{equation}\label{eq:ray}
 (S_Xr^{(d)})_k=(S_Nr)_{dk}.
\end{equation}
The amplitude and monotonicity constraints are preserved; the latter
follows by restricting the nonincreasing sequence $(1+r_m)/m$.
This is exact recursive closure, but it provides no contraction factor.

\subsection{Prime insertion in the coefficient family}
For finitely supported coefficients on squarefree labels define
\[
 F_\alpha(X,m)=
 \sum_{\substack{b:\,bm\le X\\(b,m)=1}}\mu(b)\alpha_b,\qquad
 H_\alpha(X,m)=F_\alpha(X,m)-F_\alpha(X,2m)
\]
when the labels and $m$ are odd.
Let $p$ be a fresh prime and split a new coefficient family into
$p$-free branches $\alpha_0(b)$ at label $b$ and $\alpha_1(b)$ at label
$pb$. The child branch may contain $b=1$. Direct divisor splitting gives
\begin{align}
 A_{\rm new}(n)&=A_0(n)+\ind_{p\mid n}A_1(n/p),\label{eq:insert-A}\\
 F_{\rm new}(X,m)&=F_0(X,m)-\ind_{p\nmid m}F_1(X/p,m).\label{eq:insert-F}
\end{align}
For odd $p$, odd labels, and odd $m$,
\[
 H_{\rm new}(X,m)=H_0(X,m)-\ind_{p\nmid m}H_1(X/p,m).
\]
The minus sign is $\mu(pb)=-\mu(b)$; the exclusion $p\nmid m$ is necessary.
The endpoint $X/p$ is real. Replacing $X$ by $\lfloor X\rfloor$ inside
integer inequalities is harmless; dropping cofactor exclusions or
renormalizing child masses is not.

\subsection{Signed source folding}
Let $g$ be supported on squarefree integers coprime to a fixed excluded
prime set, and let $p$ be a new prime. For $p\nmid r$ define
$(T_pg)(r)=g(r)-g(pr)$, with $g$ zero beyond its support. Pairing $r$ and
$pr$ gives
\begin{align}
 \sum_n\mu(n)g(n)&=\sum_{p\nmid r}\mu(r)(T_pg)(r),\label{eq:fold}\\
 \norm{T_pg}_1&=\norm g_1-
 2\sum_{\substack{p\nmid r\\g(r)g(pr)>0}}
                   \min(|g(r)|,|g(pr)|).\label{eq:fold-gain}
\end{align}
The function $T_pg$ remains supported on integers coprime to the excluded
primes and to $p$.
Equation~\eqref{eq:fold-gain} proves nonexpansion in $\ell^1$.
Strict contraction requires a lower bound on the displayed same-sign
overlap. It does not follow from the singular values of a two-entry
difference row.

The successful analytic argument does use a scale evolution
$\sigma\mapsto\sigma/2$, but it has not established a fixed-dyadic
contraction for these finite recursions. For a future rate improvement,
the signed prime-boundary contribution must remain coupled to composite
or cofactor sources before the final absolute values are taken.
No such stronger recursive certificate is claimed here.

\section{Relation to the literature}\label{sec:literature}

The literature comparison below was made in October 2026. Sources
marked \emph{read} were consulted in the original or in a scan. The
arithmetic bounds are classical, and the pointwise form is Wirsing's
(Section~\ref{sec:lit-wirsing}). What we compare is the method:
Section~\ref{sec:lit-mechanisms} lists the known ways of getting a
remainder from Selberg's formula and sets ours against them.

\subsection{Selberg's formula and the elementary remainder problem}

Selberg's symmetry formula \cite{selberg} is the input to the
Erd\H{o}s--Selberg proofs of the prime number theorem. As reported in \cite{wirsing1}, Selberg remarked that
his method, carried out carefully, also yields a remainder
$\psi(x)=x+O(x/(\log x)^{\alpha})$ for some $\alpha>0$. The first explicit
exponents, as listed in Wirsing's introduction \cite{wirsing1}, were
$\alpha=1/200$ (van der Corput), $\alpha=1/10$ (Kuhn \cite{kuhn}), and
$\alpha=1/6$ (Breusch \cite{breusch}, by a method starting from a formula
similar to Selberg's). Wirsing's first paper \cite{wirsing1} gave
$\alpha=3/4$ (read). Bombieri \cite{bombieri} and Wirsing's second paper
\cite{wirsing2} obtained every $\alpha$ (both read).
Later elementary work reached
$x\exp(-c(\log x)^{a})$: $a=1/7-\varepsilon$ by Diamond and Steinig, and
$a=1/6-\varepsilon$ by Lavrik--Sobirov and Sokolovski\u\i; these are cited
here from Diamond's survey \cite{diamond-survey} (read).
Selberg himself noted \cite[footnote~2]{selberg} that his method also
proves some of Beurling's theorems on generalized primes, with weaker
results than Beurling's analytic ones.

Diamond \cite[Section~8]{diamond-survey} explains the modest early exponents as
a limitation of Selberg's formula itself. Every monotone function of the
form $x/\log x+x/\log^2x+O(x/\log^cx)$ with $c>2$ satisfies the
generalized Selberg formula with remainder $O(x)$. Hence that formula
cannot detect a $\psi$-error below the scale $x/\log x$. Bombieri
overcame this with higher-weight analogues of the formula. Wirsing
instead fed each new error estimate back into an improved Selberg
formula.

\subsection{Wirsing's estimate}\label{sec:lit-wirsing}

Write $\rho(\xi)=\sum_{n\le e^\xi}(1-\Lambda(n))/n-2C$, with $C$ Euler's
constant. Wirsing \cite[Sections~2 and~4]{wirsing1} uses arithmetic
convolution identities, harmonic and divisor-sum estimates, and positivity
of $\Lambda$ to supply the two relations
\begin{equation}\label{eq:wirsing}
 \xi\rho(\xi)=\int_0^\xi\rho(\xi-\eta)\dd\rho(\eta)+O(1),\qquad
 |\rho(\xi)-\rho(\eta)|\le\xi-\eta+O(1/\eta)\quad(\eta\le\xi).
\end{equation}
The first is a Selberg-type relation for this arithmetic integrated
variable. The second is a two-sided fluctuation bound. He also notes that
$\rho(\xi)=O(\xi^{-\alpha})$ gives $\psi(x)=x+O(x/\log^\alpha x)$ with the
same $\alpha$; this follows by partial summation.

In \cite[Section~3]{wirsing1} Wirsing states that \eqref{eq:wirsing} implies
$\rho(\xi)=O(1/\xi)$, ``das bestm\"ogliche Ergebnis bei diesen
Voraussetzungen'' (the best possible result under these hypotheses).
In arithmetic terms this is the pointwise bound
\[
 \psi(x)=x+O\!\left(\frac{x}{\log x}\right),
\]
with the arithmetic input supplied in that paper and no sieve. It implies
\eqref{eq:main-energy} and is strictly stronger. This does not by itself
identify an equivalent input for an arbitrary pair $(f,\nu)$.

The distinction is visible in the input calculation
\cite[Section~4, pp.~207--208]{wirsing1}. Writing $r(x)=\rho(\log x)$,
Wirsing uses
\[
 s(x):=\sum_{m\le x}\frac{r(x/m)}m
 =\sum_{k\le x}\frac{\tau(k)-\log k-2C}{k}
 =s(\infty)+O\!\left(\frac{\log x}{\sqrt x}\right).
\]
M\"obius inversion supplies a bounded remainder in the equation for $r$;
subtracting the bounded harmonic convolution supplies the first relation
in \eqref{eq:wirsing}. These arithmetic identities are not among the base
assumptions of Theorem~\ref{thm:weak-L2}.

In \cite[Section~5]{wirsing2} the statement is made abstract, as Satz~2.
Let $\rho$ be real-valued on $[0,\infty)$, not necessarily continuous, and
assume, as $\xi\to\infty$,
\begin{itemize}
 \item $\int_0^\xi|\dd\rho|\ll\xi$;
 \item $|\rho(\xi)-\rho(\eta)|\le\gamma(\xi-\eta)+O(\eta^{-\beta})$ for
       $\eta<\xi$, with $\gamma<\sqrt2$;
 \item $\rho(\xi)=\xi^{-1}\int_0^\xi\rho(\xi-\eta)\dd\rho(\eta)
       +O(\xi^{-\beta})$ with $\beta>-1$, except at the at most countably
       many $\xi$ where the integral is undefined.
\end{itemize}
Then $\rho(\xi)\ll\xi^{-\beta}$. The example $\rho(\xi)=2\sin\xi$ shows that
$\gamma\le2$ is not enough.

Wirsing gives the proofs of the abstract theorems only in outline (``N\"aher
wollen wir auf die Beweise der beiden S\"atze nicht eingehen''). For the
arithmetic iteration he proves in full only the gain to some $\alpha'>\beta-1$,
which is all the iteration needs, and he indicates how to reach
$\alpha=\beta$ by running the induction on the whole system. His
arithmetic result $\psi(x)=x+O(x/\log^Ax)$ for every $A$ is proved in full.
It already gives $E(T)=O(1)$.

A note added in proof to \cite{wirsing2} replaces one of his lemmas by
Selberg's sieve bound for primes in short intervals. So the use of a sieve
upper density, as in Section~\ref{sec:sieve}, has a precedent; its role
there differs from its role here (Section~\ref{sec:lit-mechanisms}).

\begin{remark}[Local density is not needed on the Satz~2 route]\label{rem:ld-not-needed}
Lemma~\ref{lem:weak-density} derives
$\nu((t,t+h])\le2h+O_A(1/t)$ for bounded $h$ from the base hypotheses.
This interval bound does not by itself bound $f*f$ in Wirsing's
self-convolution equation. Theorem~\ref{thm:weak-L2} supplies the additional estimate
$f\in L^2(0,\infty)$, and Corollary~\ref{cor:wirsing-bootstrap} then
applies Satz~2 to $g=-f$. This proves $E(T)=O(T)$ for each admissible
pair without \eqref{eq:LD}. The original proof with \eqref{eq:LD}
remains a direct route with constants controlled by its hypothesis bounds.
\end{remark}

\begin{remark}[The primitive substitution does not give the bridge]
\label{rem:primitive-obstruction}
In the arithmetic case Wirsing's variable is $t-B(t)+C_0$ up to an
$O(e^{-t})$ term. For a general pair, the function $\rho(t)=t-B(t)+C_0$
need not satisfy his self-convolution equation, even under \eqref{eq:LD}.
For a fully specified example, put $\varepsilon=1/4$ and
\[
 f(t)=-e^{-t}+\varepsilon\frac{\sin\log(1+t)}{1+t},\qquad
 h=f+f',\qquad \nu=(1+h)\dd t.
\]
Here $|h|\le3\varepsilon$, so the density is positive and bounded, and
\[
 F(t)=-1+e^{-t}+\varepsilon(1-\cos\log(1+t)),\qquad
 B(t)=t+\varepsilon\left(1-\cos\log(1+t)
                   +\frac{\sin\log(1+t)}{1+t}\right).
\]
Both $f$ and $h$ are $O((1+t)^{-1})$. Hence
$f*h=O(\log(2+t)/(1+t))$ and $R=tf+F+f*h=O(1)$.
All the abstract hypotheses hold. But
\[
 \rho(t)=C_0-\varepsilon+\varepsilon\cos\log(1+t)
                  -\varepsilon\frac{\sin\log(1+t)}{1+t}
\]
is bounded, has total variation $O(\log(2+t))$, and has distinct
subsequential limits $C_0$ and $C_0-2\varepsilon$. Therefore
$\rho*d\rho=O(\log(2+t))$, whereas
$t\rho-\rho*d\rho=O(1)$ would force $\rho(t)\to0$.
No choice of $C_0$ gives that equation. This disproves the
identification, not the pointwise conclusion for $f$ or the replacement
bridge in Section~\ref{sec:wirsing-bridge}.
\end{remark}

\begin{remark}[Sharpness of the energy bound]\label{rem:energy-sharp}
Under the hypotheses of Theorem~\ref{thm:energy}, $E(T)\asymp T$ can occur.
Take, on the entire half-line,
\[
 f(t)=-e^{-t}+\frac{\sin t}{4(1+t)},\qquad
 h=f+f',\qquad \nu=(1+h)\dd t.
\]
Then $f(0)=-1$, $|h|\le3/4$, and the density is positive and bounded.
The primitive $F$ is bounded by Dirichlet's test, and the source equation
gives $B=t+f+1+F=t+O(1)$. Since $f,h=O((1+t)^{-1})$,
\[
 |(f*h)(t)|\ll\int_0^t\frac{\dd u}{(1+u)(1+t-u)}
 =\frac{2\log(1+t)}{t+2}.
\]
Thus $R=tf+F+f*h=O(1)$ and \eqref{eq:LD} holds with $q=0$.
Finally,
\[
 E(T)=\frac1{16}\int_0^T\frac{t^2\sin^2t}{(1+t)^2}\dd t+O(1)
     =\frac{T}{32}+O(\log(2+T)).
\]
This proves sharpness in the abstract class, including the stronger
density hypothesis. It does not assert sharpness for the prime-power
measure, which has additional arithmetic structure.
\end{remark}

\subsection{Mechanisms compared}\label{sec:lit-mechanisms}

The sources read use six distinct mechanisms to extract information from
Selberg's relation, or to exclude zeros on the line.
\begin{enumerate}
 \item[(K)] \emph{Sup-norm iteration in the real variable.} This is the
       route of Selberg and Erd\H{o}s, and of Kuhn \cite{kuhn} (read). An
       upper bound $|R(\eta)|\le K\eta$ on large intervals is improved by
       inserting it into the Selberg inequality. Each step loses a constant
       factor, which produces the small early exponents.
 \item[(I)] \emph{Isoperimetric estimates.} In Wirsing~I \cite{wirsing1}
       (read), and in Bang's proof described by Diamond, the Selberg
       integral is read as an area enclosed by the curve
       $(\sigma(t),\sigma(\xi-t))$, and an inequality between area and
       non-Euclidean perimeter bounds it.
 \item[(W)] \emph{Mean-square self-convolution.} Hilfssatz~1 of
       Wirsing~II \cite{wirsing2} (read) states the following. If
       $\limsup x^{-1}\int_0^x f^2=F<\infty$,
       $\limsup x^{-1}\int_0^x g^2=G<\infty$, and
       $h=x^{-1}(f*g)$ has vanishing Ces\`aro mean, then
       $\limsup x^{-1}\int_0^x h^2\le FG/2$. The constant $\tfrac12$ is
       sharp, with $\sin x$ extremal. Applied to repeated self-convolutions,
       this contracts normalized mean squares of derivatives of
       interpolants and drives Wirsing's iteration. This is not a bound
       on $\int_0^\infty|f|^2$ for our coupled source. Its application in
       Sections~3--4 of that paper already uses the self-convolution
       equation and fluctuation bound supplied by arithmetic identities.
       Bombieri's function-field note \cite{bombieri-ff} (read) closes the
       qualitative argument the same way: with $A=\limsup m^{-1}\sum r_k^2$,
       the lemma gives $A\le A^2/2$, and $A\le1$ then forces $A=0$.
 \item[(B)] \emph{Higher-weight Selberg formulas with sign-change
       reduction.} Bombieri's memoir \cite{bombieri} (read) has three
       steps. First, Amitsur's symbolic calculus, with $\log$ acting as a
       derivation, generates classes of generalized Selberg formulas of
       arbitrary weight, together with a generalized ``fundamental
       inequality''. Second, Wright's substitution $x=e^\eta$ gives
       associated functions $V_k(\eta;f)$ with
       $V_k'=V_{k-1}-V_k$. Third, the main step (his Section~VII) is a
       real-variable integral reduction after Selberg and Wright: when
       $V_0$ changes sign in an interval of length $\lambda$, its Lipschitz
       bound makes $\int|V_0|$ smaller than the majorant by a fixed factor
       $\chi<1$. This is iterated through the higher-weight formulas, with
       Beta-function bookkeeping for the convolutions. The Laplace transform
       appears only in the proof of the convolution identity
       $V_k(f*g)=\int V_mV_n$ (his Theorem~5.3), through the convolution and
       uniqueness theorems on the real axis. The memoir contains no
       frequency analysis, no doubled-frequency positivity, and no flow in
       $\sigma$.
 \item[(Z)] \emph{Nonvanishing on the line.} The 3-4-1 inequality
       $3+4\cos\theta+\cos2\theta=2(1+\cos\theta)^2\ge0$ shows that a zero
       at $1+it_0$ would force a pole at $1+2it_0$; this uses analyticity
       of $\zeta$ there. In the argument after Wiener and Geller
       \cite{wiener-geller}, presented by Diamond
       \cite[Section~7]{diamond-survey} (read), a zero makes the twisted sum
       tend to $-1$, and a short-interval upper bound rules out the
       resulting concentration of primes near bad phases. These are
       qualitative boundary arguments. The 3-4-1 argument already couples
       a frequency to its double; that coupling alone is not a new device.
 \item[(T)] \emph{Asymptotic spectral characters.} Tao \cite{tao} (read in
       summary) excludes the parity obstruction by an $|\cos|$-averaging
       gap in an asymptotic Banach algebra. It is qualitative.
\end{enumerate}

The mechanism of Sections~\ref{sec:scale}--\ref{sec:closure} has the
following components.
\begin{enumerate}
 \item[(M1)] \emph{Selberg's relation as a flow.} Along a horizontal ray
       $\sigma=ae^{-x}$, the window-normalized transform
       $\beta=\sigma b(\sigma+i\xi)$ satisfies the exact Riccati equation
       $\beta'=-\beta-\beta^2+e$ (Section~\ref{sec:scale}). Its defect is
       $O(\sigma g(\xi))$. The algebraic identity behind it,
       $\Lambda\log+\Lambda*_D\Lambda=\mu*_D\log^2$ (equivalently
       $\zeta''/\zeta=(\zeta'/\zeta)'+(\zeta'/\zeta)^2$), is classical.
       Its use as a dynamical system in $\sigma$, with stable fixed point
       $0$, unstable fixed point $-1$, and the frequency-dependent window
       $\sigma\le c/g(\xi)$ on which the defect is small, was not found in
       the sources read.
 \item[(M2)] \emph{Quantitative doubled-frequency exclusion.} (Z) closes the
       3-4-1 argument with analyticity at $1+2it_0$. Section~\ref{sec:scale}
       closes it instead with the \emph{same Riccati flow at frequency
       $2\xi$}. If $\beta(\xi)$ stays near $-1$ for a time $\log 2$, the
       variance inequality forces $\Rea\beta(2\xi)\ge4/5$, and the flow at
       $2\xi$ then decreases $\Rea\beta(2\xi)$ at a rate greater than $1$,
       which is impossible. The result $|1+\beta|\ge1/100$ holds uniformly
       in $\xi$ and in the window. No continuation and no information about
       a zeta function at $1+2it_0$ is used. The sources read have no
       analogue of this finite-time, uniform exclusion.
 \item[(M3)] \emph{Exact linearization.} The M\"obius substitution
       $W=\beta/(1+\beta)$ sends the unstable fixed point to infinity and
       turns the Riccati flow into $W'+W=e/(1+\beta)^2$ exactly. This
       yields $\int_0^{a(\xi)}|b(u+i\xi)|\dd u\le C$ uniformly in $\xi$, and
       hence a bounded variation-of-constants resolvent for the actual
       source (Section~\ref{sec:resolvent}). The fractional-linear
       substitution itself is classical. Mechanism (W) contracts
       normalized mean squares of iterated convolutions. (M3) inverts the
       Selberg operator frequency by frequency, using the actual measure
       as its own multiplier. This combined quantitative use was not
       found in the inspected sources.
 \item[(M4)] \emph{Local density as a Fourier-tail and start-line budget.}
       In (Z), and in Wirsing's note added in proof, short-interval bounds
       limit concentration of primes. Here the bound \eqref{eq:LD} instead
       controls weighted translation moduli. Through Poisson averaging with
       frequency-dependent widths, it gives an $L^2$ budget on the line
       $\sigma=a(\xi)$ (Sections~\ref{sec:tails} and~\ref{sec:start}).
       Section~\ref{sec:wirsing-bridge} uses the weaker interval bound
       derived from positivity to obtain an unweighted $L^2$ estimate.
       This yields the energy conclusion through Wirsing's theorem.
       The direct weighted-tail proof still uses \eqref{eq:LD}; the new
       route does not assert that the same weighted estimates hold with
       an arbitrary $O(1/t)$ density error.
 \item[(M5)] \emph{$L^2$ charging of the forcing.} The Selberg remainder
       $R$ enters only through $\int|P|^2$, by Tonelli
       (Section~\ref{sec:closure}); a pointwise charge would lose a
       logarithm. This is a standard device.
\end{enumerate}

The proof contribution is the combination of (M1)--(M3) with
the tail and source estimates: a quantitative, frequency-resolved,
continuation-free treatment of Selberg's relation in the half-plane of
absolute convergence. The comparison distinguishes this combination
from the inspected real-variable and qualitative boundary arguments;
it does not establish historical priority for the combination or its
individual ingredients.

\emph{Scope of the search.} The prime number theorem literature is too large
to search exhaustively. The search covered Diamond's survey of elementary
methods \cite{diamond-survey}, which catalogs the Tauberian routes from
Selberg's formula, together with the primary sources for each distinct
remainder mechanism compared here (Kuhn, Wirsing~I and~II, Bombieri).
We inspected all of Wirsing~I, pp.~205--214, and
Wirsing~II, Sections~1--5, pp.~1--15, and found no matched supply of global
$L^2$ control under precisely our base coupled-source assumptions.
This is a bounded source finding, not a proof that no classical derivation
exists. For the Bombieri memoir we read pages 1--6, 21--34, and 39--46
of the scan we used, not every page. The search did not cover
every derivative paper. Ingham's review \cite{ingham} was not consulted;
Wright \cite{wright} and Amitsur \cite{amitsur} are known to us through
Bombieri's use of them in (B). No conclusion about their unexamined
contents is drawn.

\subsection{Other related work}

\begin{itemize}
 \item Tao \cite{tao} (read in summary) uses Selberg's formula and
       positivity of the log-prime measure in an asymptotic Banach-algebra
       proof of the prime number theorem. It is qualitative and gives no
       rate.
 \item Saidak \cite{saidak} (read) proves the reverse implication: a PNT
       remainder gives a sharpened Selberg formula. His Remark~2.5 notes
       that the elementary-remainder literature did not sharpen Selberg's
       lemma itself.
 \item Hildebrand \cite{hildebrand} (first page read) proves $M(x)=o(x)$
       without Selberg-type formulas, using a large-sieve inequality
       together with Mertens' estimates and a sieve upper bound for primes
       in short intervals.
 \item Debruyne and Vindas \cite{dv} develop Laplace-Tauberian remainder
       theorems for Beurling generalized primes.
 \item Kahane \cite{kahane} proves an $L^2$ prime number theorem for
       Beurling systems, with an $L^2$ hypothesis on the integer-counting
       error and an analytic proof. It is not a statement about abstract
       Selberg-type relations.
 \item McNamara \cite{mcnamara} gives a dynamical proof of the prime
       number theorem that also uses Selberg's formula, through an
       entropy-decrement argument over real averages. It uses neither
       the Laplace transform nor a flow in $\sigma$, and gives no rate.
 \item The monograph of Diamond and Zhang \cite{diamond-zhang} treats
       Beurling systems and $L^2$ prime number theorems (Kahane), using both
       analytic and elementary methods. It is a compilation and was not
       consulted.
\end{itemize}

\emph{The author's earlier work.} The certificate of
Sections~\ref{sec:certificate}--\ref{sec:transfer} continues the author's
work on sparse matrix representations of the Mertens function
\cite{kline2019,kline2020,kline2020b,kline-sv,kline-ss}. The matrix of
\cite{kline2019} has determinant $M(n)$, and its rows annihilate the
M\"obius vector. Section~7.1 of \cite{kline-ss} notes that any vector $z$
in the row space gives $M(n)=\mu^{\mathsf T}(\mathbf 1-H^{\mathsf T}z)$, and
names an explicit choice of coefficients with a uniform residual bound as
the missing ingredient. Proposition~\ref{prop:certificate} uses a
different family of relations, $\mu(bm)=\mu(b)\mu(m)$ over divisor labels
$b$, with an $\ell^1$ residual, and Section~\ref{sec:weights} supplies
explicit weights. The operators $S_N$, $Z_N$, and $Q_N$ of
Section~\ref{sec:finite} do not occur in those works.

\subsection{Contribution}

\begin{enumerate}
 \item The method (M1)--(M3). Selberg's formula is treated as a Riccati
       equation along horizontal lines. Positivity at the doubled frequency
       keeps the solution away from its unstable fixed point, uniformly and
       in finite time. A M\"obius substitution removes the quadratic term,
       leaving a linear equation with a small solution-dependent forcing,
       and gives a uniform resolvent. With (M4)--(M5), this is a complete proof
       of the energy bound under \eqref{eq:LD}, without analytic
       continuation.
 \item The $L^2$ estimate of Section~\ref{sec:wirsing-bridge}, which needs
       no density hypothesis. It supplies the input that Wirsing's abstract
       theorem requires for a general positive measure. Wirsing obtained
       that input from arithmetic identities, which a general measure does
       not have (Remark~\ref{rem:primitive-obstruction}).
 \item The explicit passage from the energy bound to the sparse Mertens
       certificate (Sections~\ref{sec:certificate}--\ref{sec:transfer}).
\end{enumerate}
We did not find (1) or (2) in the sources read. That search was limited, so
priority is not settled.

\section{Position of the result and remaining questions}\label{sec:position}

The divisor identities, the prime and semiprime weights, Selberg's formula,
and the quadratic sieve are classical. So are the arithmetic bounds, and
the abstract theorem used in Corollary~\ref{cor:wirsing-bootstrap} is
Wirsing's. This paper contributes the Riccati-flow method (M1)--(M3), the
estimates it gives for the source, and the explicit chain to the sparse
certificate.

Status:
\begin{itemize}
 \item Proved here: the energy theorem, the density bound for primes, the
       cofactor transfer, and \eqref{eq:arith-conclusions}.
 \item Proved here without \eqref{eq:LD}: $f\in L^2(0,\infty)$
       (Theorem~\ref{thm:weak-L2}).
 \item Proved here, assuming Wirsing's Satz~2: the energy bound and
       $f(t)=O(1/t)$ without \eqref{eq:LD}. The constant may depend on
       the particular pair $(f,\nu)$.
 \item Sharp: the energy bound cannot be improved in the abstract setting
       (Remark~\ref{rem:energy-sharp}).
 \item Open: the general constrained inverse \eqref{eq:constrained-open},
       the spectral tail \eqref{eq:svd-tail-open}, and a uniformly
       contracting finite recursion.
 \item Limit: the current certificate cannot go below the floor in
       Remark~\ref{rem:floor}, so the method gives no better rate as it
       stands.
\end{itemize}
The next step is a signed certificate that can go below that floor,
together with the estimate it would need. Larger numerical sweeps, or
rewriting the same finite identity, will not supply that estimate.

\appendix
\section{A circular source estimate and its transfer limitation}
\label{app:circle}

The following result records a direct singular-mode mechanism.
It is not needed for the exponential-window proof, but explains why the
source must stay linked to the measure.

\begin{proposition}\label{prop:circle}
Let $\eta$ be any probability measure on the circle $\R/T\mathbb Z$,
$T>0$. Let $q$ be the mean-zero periodic solution of
\[
 dq+q\dd t=T\eta-\dd t.
\]
Then
\begin{equation}\label{eq:circle-bound}
 \norm q_2\le\frac{15}{2}\norm{q+\eta*q}_2.
\end{equation}
\end{proposition}
\begin{proof}
Use normalized Fourier coefficients, and put
\[
 \omega_j=2\pi j/T,\qquad
 z_j=\int e^{-i\omega_jt}\eta(\mathrm dt),\qquad
 w_j=(1+\omega_j^2)^{-1}.
\]
For $j\ne0$, $q_j=z_j/(1+i\omega_j)$; $q_0=0$.
These coefficients define an $L^2$ solution since $\sum_jw_j<\infty$.
Set $e_j=w_j|z_j|^2$ and $d_j=w_j|z_j(1+z_j)|^2$.
Call $j$ bad if $|1+z_j|<1/6$. Good indices satisfy $e_j\le36d_j$.
At a bad index, $\Rea z_j<-5/6$, and positivity gives
\[
 \Rea z_{2j}\ge2(\Rea z_j)^2-1>7/18>1/3.
\]
Therefore $|z_{2j}|\ge1/3$, $|1+z_{2j}|\ge4/3$, and
$w_{2j}\ge w_j/4$, so $d_{2j}\ge(4/81)e_j$.
The map $j\mapsto2j$ is injective. For every integer $J\ge1$,
\[
 \sum_{0<|j|\le J}e_j
 \le36\sum_{0<|j|\le J}d_j+
       \frac{81}{4}\sum_{0<|j|\le2J}d_j
 \le\frac{225}{4}\sum_{0<|j|\le2J}d_j.
\]
This explicitly pays for any overlap between directly controlled and
charged modes. Let $J\to\infty$ and apply Parseval to obtain
\eqref{eq:circle-bound}.
\end{proof}

If $\eta$ is a point mass at $T/2$, the operator $I+\eta*$ annihilates
every odd Fourier mode. Thus \eqref{eq:circle-bound} is not an all-vector
singular-value bound. For the generated source, the even companion
modes control the bad odd modes in aggregate.
If $v=q+\varepsilon$ with $\norm\varepsilon_2\le e_0$, then
$\norm{I+\eta*}\le2$ and \eqref{eq:circle-bound} give
\[
 \norm v_2\le\frac{15}{2}\norm{(I+\eta*)v}_2+16e_0.
\]

For the arithmetic measure, take large $T$ away from its atoms and set
\[
 B_T=B(T),\quad \eta_T=\nu|_{[0,T]}/B_T,\quad
 c_T=T/B_T,\quad b_T=\frac{c_T(1+f(T))}{1-e^{-T}}.
\]
On $0\le t<T$, define
\begin{equation}\label{eq:circle-lift}
 q_T(t)=c_Tf(t)+(c_T-1)+b_Te^{-t}.
\end{equation}
Its endpoint traces agree, and
$dq_T+q_T\dd t=T\eta_T-\dd t$; integrating gives zero mean.
Also $c_T=1+O(1/T)$ and $b_T=O(1)$, so
\begin{equation}\label{eq:circle-lift-cost}
 \int_0^Tt^2|q_T(t)-f(t)|^2\dd t\ll T.
\end{equation}
For the projection $\Pi_J$ onto $0<|j|\le J$, the bound $|z_j|\le1$ gives
\[
 \int_0^T|q_T-\Pi_Jq_T|^2\dd t\le\frac{T^3}{2\pi^2J},\qquad
 \int_0^Tt^2|q_T-\Pi_Jq_T|^2\dd t\le\frac{T^5}{2\pi^2J}.
\]
These are analytic tail bounds, not finite-data extrapolations.

The circular residual still has a wrap term. Let $K_T=\eta_T*$,
$e_T(t)=c_T-1+b_Te^{-t}$, and
\[
 W_T(t)=\int_{(t,T]}f(T+t-u)\nu(\mathrm du).
\]
Splitting the circular convolution gives, almost everywhere,
\begin{equation}\label{eq:wrap}
 (I+K_T)q_T
 =c_T\left[\left(1-\frac t{B_T}\right)f(t)
                  +\frac{W_T(t)+R(t)}{B_T}\right]+e_T+K_Te_T.
\end{equation}
Replacing this residual by $R/B_T$ would omit both the early source and
the wrap term. Proposition~\ref{prop:circle} alone therefore does not close
the weighted arithmetic target. Sections~\ref{sec:scale}--\ref{sec:closure}
give a complete alternative using exponential windows and retaining
the initial-source term explicitly.

\section*{AI assistance}

The arguments in this paper were produced with the assistance of AI models
and checked by separate AI proof audits. Agreement between AI audits is not
independent human verification, and the paper has not been peer reviewed.

\begin{thebibliography}{99}
\bibitem{bombieri}
E. Bombieri,
\emph{Maggiorazione del resto nel ``Primzahlsatz'' col metodo di
Erd\H{o}s--Selberg}, Ist. Lombardo Accad. Sci. Lett. Rend. A \textbf{96}
(1962), 343--350; and \emph{Sulle formule di A. Selberg generalizzate per
classi di funzioni aritmetiche e le applicazioni al problema del resto nel
``Primzahlsatz''}, Riv. Mat. Univ. Parma (2) \textbf{3} (1962), 393--440.

\bibitem{amitsur}
S. A. Amitsur,
\emph{Arithmetic linear transformations and abstract prime number theorems},
Canad. J. Math. \textbf{13} (1961), 83--109.

\bibitem{ingham}
A. E. Ingham,
review of A. Selberg, \emph{An elementary proof of the prime-number
theorem}, and P. Erd\H{o}s, \emph{On a new method in elementary number
theory which leads to an elementary proof of the prime number theorem};
reprinted in \emph{Reviews in Number Theory as printed in Mathematical
Reviews 1940--1972}, vol.~4, N20-3, 191--193, Amer. Math. Soc., 1974.

\bibitem{wright}
E. M. Wright,
\emph{The elementary proof of the prime number theorem},
Proc. Roy. Soc. Edinburgh Sect. A \textbf{63} (1952), 257--267.

\bibitem{bombieri-ff}
E. Bombieri,
\emph{Sull'analogo della formula di Selberg nei corpi di funzioni},
Atti Accad. Naz. Lincei Rend. Cl. Sci. Fis. Mat. Natur. (8) \textbf{35}
(1963), 252--257.

\bibitem{breusch}
R. Breusch,
\emph{An elementary proof of the prime number theorem with remainder term},
Pacific J. Math. \textbf{10} (1960), 487--497.

\bibitem{diamond-survey}
H. G. Diamond,
\emph{Elementary methods in the study of the distribution of prime numbers},
Bull. Amer. Math. Soc. (N.S.) \textbf{7} (1982), 553--589.

\bibitem{diamond-zhang}
H. G. Diamond and W.-B. Zhang,
\emph{Beurling Generalized Numbers},
Math. Surveys Monogr. \textbf{213}, Amer. Math. Soc., 2016.

\bibitem{hildebrand}
A. Hildebrand,
\emph{The prime number theorem via the large sieve},
Mathematika \textbf{33} (1986), 23--30.

\bibitem{kuhn}
P. Kuhn,
\emph{Eine Verbesserung des Restgliedes beim elementaren Beweis des
Primzahlsatzes}, Math. Scand. \textbf{3} (1955), 75--89.

\bibitem{wiener-geller}
N. Wiener and L. Geller,
\emph{Some prime-number consequences of the Ikehara theorem},
Acta Sci. Math. (Szeged) \textbf{12B} (1950), 25--28.

\bibitem{wirsing1}
E. Wirsing,
\emph{Elementare Beweise des Primzahlsatzes mit Restglied.~I},
J. Reine Angew. Math. \textbf{211} (1962), 205--214.

\bibitem{wirsing2}
E. Wirsing,
\emph{Elementare Beweise des Primzahlsatzes mit Restglied.~II},
J. Reine Angew. Math. \textbf{214/215} (1964), 1--18.

\bibitem{selberg}
A. Selberg,
\emph{An elementary proof of the prime-number theorem},
Annals of Mathematics (2) \textbf{50} (1949), 305--313.

\bibitem{tao}
T. Tao,
\emph{A Banach algebra proof of the prime number theorem},
25 October 2014.
\href{https://terrytao.wordpress.com/2014/10/25/a-banach-algebra-proof-of-the-prime-number-theorem/}
{Online exposition}.

\bibitem{saidak}
F. Saidak,
\emph{On the prime number lemma of Selberg},
Mathematica Scandinavica \textbf{103} (2008), 5--10.
\href{https://doi.org/10.7146/math.scand.a-15065}
{doi:10.7146/math.scand.a-15065}.
An author preprint is available at
\url{https://uncg.edu/~f_saidak/MS2007.pdf}.

\bibitem{dv}
G. Debruyne and J. Vindas,
\emph{On General Prime Number Theorems with Remainder},
arXiv:1601.05324, version 2, 2017.
\url{https://arxiv.org/abs/1601.05324}.
Published in \emph{Generalized Functions and Fourier Analysis}, Oper. Theory
Adv. Appl. \textbf{260}, Birkh\"auser, 2017, 79--94,
\href{https://doi.org/10.1007/978-3-319-51911-1_6}{doi:10.1007/978-3-319-51911-1\_6}.

\bibitem{selberg-sieve}
A. Selberg,
\emph{On an elementary method in the theory of primes},
Norske Vid. Selsk. Forh. (Trondheim) \textbf{19} (1947), 64--67.

\bibitem{mv}
H. L. Montgomery and R. C. Vaughan,
\emph{The large sieve},
Mathematika \textbf{20} (1973), 119--134,
\href{https://doi.org/10.1112/S0025579300004708}{doi:10.1112/S0025579300004708}.

\bibitem{kahane}
J.-P. Kahane,
\emph{Sur les nombres premiers g\'en\'eralis\'es de Beurling. Preuve d'une
conjecture de Bateman et Diamond},
J. Th\'eor. Nombres Bordeaux \textbf{9} (1997), 251--266,
\href{https://doi.org/10.5802/jtnb.201}{doi:10.5802/jtnb.201}.

\bibitem{mcnamara}
R. McNamara,
\emph{A dynamical proof of the prime number theorem},
Hardy--Ramanujan J. \textbf{44} (2021),
\href{https://doi.org/10.46298/hrj.2022.8924}{doi:10.46298/hrj.2022.8924}.

\bibitem{kline2019}
J. Kline,
\emph{A sparser matrix representation of the Mertens function},
Linear Algebra Appl. \textbf{581} (2019), 354--366,
\href{https://doi.org/10.1016/j.laa.2019.07.021}{doi:10.1016/j.laa.2019.07.021}.

\bibitem{kline2020}
J. Kline,
\emph{On the eigenstructure of sparse matrices related to the prime number theorem},
Linear Algebra Appl. \textbf{584} (2020), 409--430,
\href{https://doi.org/10.1016/j.laa.2019.09.022}{doi:10.1016/j.laa.2019.09.022}.

\bibitem{kline2020b}
J. Kline,
\emph{Bordered Hermitian matrices and sums of the M\"obius function},
Linear Algebra Appl. \textbf{588} (2020), 224--237,
\href{https://doi.org/10.1016/j.laa.2019.12.004}{doi:10.1016/j.laa.2019.12.004}.

\bibitem{kline-sv}
J. Kline,
\emph{The smallest singular value of a sparse Mertens matrix}, research
release v0.1.0, 2026,
\href{https://doi.org/10.5281/zenodo.21774716}{doi:10.5281/zenodo.21774716}.

\bibitem{kline-ss}
J. Kline,
\emph{The singular spectrum of a sparse Mertens matrix}, research release
v0.1.0, 2026,
\href{https://doi.org/10.5281/zenodo.23004980}{doi:10.5281/zenodo.23004980}.
\end{thebibliography}
\end{document}
